\documentclass{article}

\usepackage{packages}
\usepackage{environments}
\usepackage{commands}

\title{Написание кода на C++}
\author{Дима Трушин}
\date{2024 --- 2025}


\begin{document}


\maketitle
\tableofcontents

\newpage

\section{Память}

Я хочу начать разговор про память с неожиданной стороны.
В отличие от Си в C++ более богат на сущности.
И это выражается не только синтаксически в виде написания классов, но и в идейной работе языка и компилятора.
В частности в C++ есть понятие объекта, который может занимать какую-то память.
У этого объекта есть всякие свойства вроде времени жизни.
С точки зрения Си у вас любой кусок памяти нужного размера и правильно выровненный можно считать структурой заполненной мусором.
В C++ так уже нельзя.
Тут надо думать про сырую память не занимаемую ничем и про занятую объектом память.
В частности обращение к сырой памяти, если там нет объекта или не тот объект вызывает UB.
Давайте разбираться что это такое.

\subsection{Объекты}

Язык C++ использует понятие объекта для описания данных, с которыми происходят манипуляции.
Это могут быть данные встроеного типа (такие как \verb"int" или \verb"double") или класса лежащие в переменной или временные объект, объект возвращаемый из функции и т.д.
Ментальная модель такая: переменные -- это стаканы для данных, данные внутри -- это объекты.
Еще бывают временные стаканы, в них лежат временные объекты.
Самое важное для нас -- у объектов в C++ value semantics, что по простому означает, что если у вас есть два разных объекта, то никакие изменения одного не могут повлечь изменение другого и наоборот.
Это создает логическую локальность для данных в коде.
Тут стоит сказать, что в языке есть объекты с ограниченным reference semantics и можно самостоятельно имплементировать reference semantics полноценно.
Более подробно семантику языка я обсуждал в разделе~\ref{}.

Важно понимать, что любой объект в C++ обладает следующими свойствами:
\begin{enumerate}
\item Размер объекта (раздел~\ref{section::ObjSizeAlignment}).
Размер определяет сколько места в памяти занимает объект.
Размер определяется встроенной функцией \verb"sizeof".


\item Требования по выравниванию или \verb"alignment" (раздел~\ref{section::ObjSizeAlignment}).
Требования к выравниванию определяют по какому адресу могут находиться начало и конец объекта.
Требования к выравниванию определяется встроенной функцией \verb"alignof".

\item Storage duration -- время жизни выделенной под объект памяти (раздел~\ref{section::StorageDuration}).
Бывает: \verb"automatic", \verb"static", \verb"dynamic", \verb"thread-local".

\item \verb"Lifetime" или время жизни объекта (раздел~\ref{section::Lifetime}).
Логическое время жизни объекта, грубо говоря, с момента окончания работы конструктора, до момента вызова деструктора.
Время жизни ограничено storage duration или условиями на временный объект.

\item Тип.

\item \verb"Value" или значение (раздел~\ref{section::ObjValue}).
Значение -- это конкретное состояние объекта.
Значение у объекта может быть неопределенным, например для не инициализированных переменных встроенного типа.

\item Имя -- опционально.
Обычно это имя переменной, в которой хранится объект.
\end{enumerate}
Для понимания того, как можно и нужно работать с памятью полезно немного понимать про модель объектов в C++.
Ниже я буду по очереди разбирать все эти понятия с примерами и объяснением для чего нужно знать подобные вещи.

\paragraph{Что является объектами}

% TO DO
% Нужны примеры объектов
TO DO

\paragraph{Что не является объектами}

Чтобы хоть немного сориентироваться, проговорим явно, что следующие сущности не являются объектами:
\begin{enumerate}
\item \verb"Value" (значения).
По сути объект представлен всеми данными в памяти от адреса его начала и до адреса его конца.
А \verb"value" отвечает за его логическое представление.
Например, у объекта могут быть padding-ги из-за выравнивания (см.~раздел~\ref{section::ObjSizeAlignment}), а потому биты входящие в них никак не влияют на значение объекта.
Так же бывает, что какие-то состояния у объекта недопустимы.
Например ссылка это по сути указатель с красивым синтаксисом.
Но ссылка не может быть нулевой.

\item Ссылки не относятся к объектам.
Это разумно, так как у ссылок совсем другая семантика.
Первое равенство для ссылок -- это бинд ссылки на конкретный объект, а все остальные использования ссылки ничем не отличимы от использования имени объекта напрямую.

\item Функции не относятся к объектам.
Функции лежат в памяти недоступной для программиста.
Вот указатели на функции -- это объекты.

\item Элементы \verb"enum"-ов не являются объектами.
То есть если вы пишите объявление
\begin{cppcode}[numbers = none]
enum class E : unsigned char {
  First, Second, Last
};
\end{cppcode}
То элементы \verb"E::First", \verb"E::Second", \verb"E::Last" не являются объектами.

\item Типы не являются объектами.
То есть встроенные типы вида \verb"int" и \verb"double" -- это не объекты с точки зрения языка.
Так же любые классы, которые мы объявляем
\begin{cppcode}[numbers = none]
struct A  {};

union B {
  int x;
  short y;
};
\end{cppcode}
Классы \verb"A" и \verb"B" -- это не объекты с точки зрения языка.%
\footnote{Обратите внимание, что \verb"union" это тоже класс с точки зрения языка.}

\item Не статический член класса не являются объектами.
\begin{cppcode}[numbers = none]
struct A  {
  void f() const {}
  int x;
};
\end{cppcode}
В этом случае \verb"A::x" и \verb"A::f" -- это не объекты с точки зрения языка.
Тут логика в том, что самих по себе объектов \verb"A::x" не существует, они есть только лишь когда есть содержащий их объект.
А \verb"A::f" не существует, потому что это на самом деле просто функция (хотя программисту доступны лишь указатели на них), а функции мы не считаем объектами.

\item Шаблоны не являются объектами.
\begin{cppcode}[numbers = none]
template<class>
struct A {};

template<class>
void f() {}

template<class>
constexpr bool val = {};
\end{cppcode}
Все выше перечисленное не является объектами.
Это <<чертежи>> для создания типов, функций и переменных.
Последние (переменные) уже являются именами объектов, которые они содержат.

\item Специализация шаблона класса или функции не является объектом.
\begin{cppcode}[numbers = none]
template<class>
void f() {}

template<>
void f<int>() {}

template<class>
struct A {};

template<>
struct A<int> {};
\end{cppcode}
Все выше перечисленное не объекты.
Специализация шаблона функции дает просто функцию, а они не объекты.
А специализация шаблона класса, дает класс, который тоже не объект.

\item Пространство имен тоже не относится к объектам.
\begin{cppcode}[numbers = none]
namespace detail {
  struct A {};
}
\end{cppcode}
В данном случае \verb"detail" не является объектом.

\item Parameter pack в шаблонах не является объектом.
\begin{cppcode}[numbers = none]
template<class... Args>
struct A {};
\end{cppcode}
Конструкция \verb"Args" называется parameter pack и не является объектом.

\item Неявный аргумент всех методов класса \verb"this" тоже не является объектом.
Потому для него не доступны все те же свойства, что и обычным указателям (которые объектами являются).
\end{enumerate}

\subsection{Размер и выравнивание}
\label{section::ObjSizeAlignment}

Теперь обсудим самые базовые вещи -- размер и выравнивание данных в языке.

\subsubsection{Размер и выравнивание базовых типов}

В языке C++ память, с которой работает программист представляет из себя непрерывный массив из байт.
Представлять его себе можно как \verb"std::vector".
\[
\noteblock{8}{}{\texttt{address = 0}}{4}
\noteblock{8}{}{\texttt{address = 8}}{4}
\noteblock{8}{}{\texttt{address = 16}}{4}
\note{4}{\texttt{address = 24}}
\lastblock{}
\]
Не все адреса одинаково доступны программисту, но это для нас сейчас не важно.
Любой объект занимает в памяти непрерывный сегмент.
Размер этого сегмента -- размер объекта.
Начнем со встроенных типов данных языка:
\[
\noteblock{4}{}{\texttt{0}}{1}
\datablock{4}{\cellcolor{Orange!40}}{\texttt{4 bytes}}{\texttt{int}}{2}
\noteblock{2}{}{\texttt{8}}{1}
\datablock{2}{\cellcolor{OliveGreen!40}}{\texttt{2 bytes}}{\texttt{short}}{2}
\emptyblock{2}{}
\datablock{1}{\cellcolor{Magenta!40}}{\texttt{1 bytes}}{\texttt{char}}{1}
\emptyblock{1}{}
\note{1}{\texttt{16}}
\lastblock{}
\]
Здесь я полагаю, что размер \verb"char" -- $1$ байт, размер \verb"short" -- $2$ байта и размер \verb"int" -- $4$ байта.
Кроме размера у каждого типа есть требования к выравниванию -- \verb"alignment".
Выравнивание означает, что данные этого типа не могут начинаться и заканчиваться в памяти где угодно.%
\footnote{Причина почему данные не могут заканчиваться где угодно в том, что если мы хотим создать массив объектов, то они в памяти будут лежать друг за другом.
А это значит что там, где заканчивается один объект -- начинается другой.
А потому надо выровнять концовку объекта тоже.
Кроме того, стэк растет в противоположную сторону и для корректного помещения объекта на стек также удобно, когда оба конца выровнены.}
Для типа \verb"T" выравнивание равно $d$, если адрес объекта типа \verb"T" должен делиться на $d$.
Для встроенных типов данных \verb"alignment" совпадает с размером.
В частности в нашем примере верно следующее:
\begin{enumerate}
\item Данные типа \verb"char" могут располагаться где угодно.
\[
\noteblock{4}{}{\texttt{0}}{1}
\datablock{1}{\cellcolor{Magenta!40}}{\texttt{1 bytes}}{\texttt{char}}{1}
\emptyblock{2}{}
\datablock{1}{\cellcolor{Magenta!40}}{\texttt{1 bytes}}{\texttt{char}}{1}
\emptyblock{3}{}
\datablock{1}{\cellcolor{Magenta!40}}{\texttt{1 bytes}}{\texttt{char}}{1}
\emptyblock{4}{}
\note{1}{\texttt{16}}
\lastblock{}
\]

\item Данные типа \verb"short" обязаны начинаться с четного адреса.
\[
\noteblock{2}{}{\texttt{0}}{1}
\datablock{2}{\cellcolor{OliveGreen!40}}{\texttt{2 bytes}}{\texttt{short}}{2}
\emptyblock{4}{}
\datablock{2}{\cellcolor{OliveGreen!40}}{\texttt{2 bytes}}{\texttt{short}}{2}
\emptyblock{6}{}
\note{1}{\texttt{16}}
\lastblock{}
\]
Вот такое расположение типа \verb"short" в памяти недопустимо
\[
\noteblock{1}{}{\texttt{0}}{1}
\datablock{2}{\cellcolor{red!40}}{\texttt{2 bytes}}{\texttt{short}}{2}
\emptyblock{4}{}
\datablock{2}{\cellcolor{red!40}}{\texttt{2 bytes}}{\texttt{short}}{2}
\emptyblock{7}{}
\note{1}{\texttt{16}}
\lastblock{}
\]

\item данные типа \verb"int" обязаны начинаться с адреса делящегося на $4$.
\[
\noteblock{4}{}{\texttt{0}}{1}
\datablock{4}{\cellcolor{Orange!40}}{\texttt{4 bytes}}{\texttt{int}}{2}
\datablock{4}{\cellcolor{Orange!40}}{\texttt{4 bytes}}{\texttt{int}}{2}
\emptyblock{4}{}
\note{1}{\texttt{16}}
\lastblock{}
\]
Вот такое расположение типа \verb"int" в памяти недопустимо
\[
\noteblock{2}{}{\texttt{0}}{1}
\datablock{4}{\cellcolor{red!40}}{\texttt{4 bytes}}{\texttt{int}}{2}
\emptyblock{1}{}
\datablock{4}{\cellcolor{red!40}}{\texttt{4 bytes}}{\texttt{int}}{2}
\emptyblock{5}{}
\note{1}{\texttt{16}}
\lastblock{}
\]
\end{enumerate}
Подобные ограничения связаны с тем, как работает железо.
Процессор не может просто так прочитать данные из любого сегмента памяти.
При попытке обратиться к неверно выровненным данным, вы получаете UB.

\subsubsection{Простые структуры}

Теперь обсудим, как вычисляется размер и выравнивание классов.
Начнем с простого случая -- обычная структура, без базовых классов, без виртуальных методов.
Рассмотрим следующие два класса:
\begin{center}
\begin{tabular}{ll}
{
\begin{minipage}[\baselineskip]{6cm}
\begin{cppcode}[numbers = none]
struct A {
  int x;
  short y;
  short z;
};
\end{cppcode}
\end{minipage}
}&{
\begin{minipage}[\baselineskip]{6cm}
\begin{cppcode}[numbers = none]
struct B {
  short y;
  int x;
  short z;
};
\end{cppcode}
\end{minipage}
}
\end{tabular}
\end{center}
Данные в памяти должны быть размещены с учетом требований по размеру и выравниванию ко всем полям.
Для выполнения этих требований компилятор использует следующие правила:
\begin{enumerate}
\item Размер структуры должен быть кратен ее выравниванию.

\item Выравнивание класса равно максимуму среди выравниваний по всем полям класса.

\item Размер вычисляется последовательным расположением данных в памяти так, чтобы каждый член был правильно выровнен.
Можно считать, что мы располагаем данные начиная с адреса $0$.
\end{enumerate}
У обоих классов выравнивание будет равно $4$, так как это максимум между выравниванием для \verb"int" и \verb"short".
А вот с размером будет небольшая разница из-за того, что мы должны положить в память данные в том же порядке, в каком они идут в классе.
Вот что будет в случае класса \verb"A".
\[
\noteunder{8}{\texttt{A}}
\noteblock{4}{\cellcolor{Orange!40}}{\texttt{x}}{1}
\noteblock{2}{\cellcolor{OliveGreen!40}}{\texttt{y}}{1}
\noteblock{2}{\cellcolor{OliveGreen!40}}{\texttt{z}}{1}
\]
Адреса будут $0$, $4$, $6$.
В этом случае размер класса будет $8$ байт.
Однако, если мы посмотрим на класс \verb"B", то расположение в памяти будет следующее:
\[
\noteunder{12}{\texttt{B}}
\noteblock{2}{\cellcolor{OliveGreen!40}}{\texttt{y}}{1}
\noteblock{2}{}{\texttt{pad}}{2}
\noteblock{4}{\cellcolor{Orange!40}}{\texttt{x}}{1}
\noteblock{2}{\cellcolor{OliveGreen!40}}{\texttt{z}}{1}
\noteblock{2}{}{\texttt{pad}}{2}
\]
После расположения в памяти поля $y$ типа \verb"short", следующий свободный адрес будет $2$ и нам надо расположить поле \verb"x" типа \verb"int" следом.
Однако, оно может располагаться только по адресу кратному $4$.
В этом случае мы добавляем $2$ пустых байта (padding) и складываем поле \verb"x" по адресу $4$.
После чего по адресу $8$ складываем еще одно поле \verb"z" типа \verb"short".
И надо в конце добавить пустые байты, чтобы размер делился на alignment.
Это еще $2$ пустых байта.
Последние байты нужны для того, чтобы при помещении элементов класса в массив (они в памяти располагаются друг за другом), следующий элемент располагался по допустимому относительно выравнивания адресу.

\subsubsection{Наличие базовых классов}

Теперь посмотрим, что будет, если у нас есть какое-то количество базовых классов.
\begin{cppcode}
struct C : A {
  bool f;
  B b;
  double d;
};
\end{cppcode}
Выравнивание все еще считается как максимум выравниваний по всем базовым классам и полям класса.
В данном случае максимум будет $8$ из-за наличия \verb"double".
А размер считается по той же самой процедуре, только сначала мы должны в памяти расположить базовые классы, в порядке слева направо, а потом уже поля класса сверху вниз.%
\footnote{Просто трактуйте базовый класс как переменную класса, у которой нет имени, и к ней надо обращаться по ее типу.}
Получим следующее:
\[
\noteunder{8}{\texttt{A}}
\noteblock{4}{\cellcolor{Orange!40}}{\texttt{x}}{1}
\noteblock{2}{\cellcolor{OliveGreen!40}}{\texttt{y}}{1}
\noteblock{2}{\cellcolor{OliveGreen!40}}{\texttt{z}}{1}
\noteunder{1}{\texttt{bool}}
\noteblock{1}{\cellcolor{TealBlue!40}}{\texttt{f}}{1}
\noteblock{3}{}{\texttt{pad}}{2}
\noteunder{12}{\texttt{B}}
\noteblock{2}{\cellcolor{OliveGreen!40}}{\texttt{y}}{1}
\noteblock{2}{}{\texttt{pad}}{2}
\noteblock{4}{\cellcolor{Orange!40}}{\texttt{x}}{1}
\noteblock{2}{\cellcolor{OliveGreen!40}}{\texttt{z}}{1}
\noteblock{2}{}{\texttt{pad}}{2}
\noteunder{8}{\texttt{double}}
\noteblock{8}{\cellcolor{Blue!40}}{\texttt{d}}{2}
\]
В итоге размер будет $32$ байта.

\subsubsection{Наличие пустого базового класса}

Стоит сказать про наличие оптимизации пустого класса.
По умолчанию, если класс не содержит данных, то стандарт требует, чтобы его объект был ненулевого размера.
То есть для класса
\begin{cppcode}
struct A {
  void f() const {}
};
\end{cppcode}
размер будет $1$ байт и выравнивание $1$.
Связано это с тем, что модель С++ изначально не допускала расположения по одному адресу нескольких объектов (хотя сейчас допускает), что могло произойти в следующем случае
\begin{cppcode}
struct B {
  A a1;
  A a2;
};
\end{cppcode}
В этом случае \verb"B" будет занимать $2$ байта и выравнивание будет $1$.
Однако, если мы унаследуемся от пустого класса, то он не будет занимать места.
Например в случае
\begin{cppcode}
struct A {
  void f() const {}
};

struct B{
  int x;
};

struct C {
  short y;
};

struct D : B, A, C {
  char z;
};
\end{cppcode}
выравнивание будет $4$, размер $8$ и данные будут лежать в памяти следующим образом
\[
\noteunder{4}{\texttt{B}}
\noteblock{4}{\cellcolor{Orange!40}}{\texttt{x}}{1}
\noteunder{2}{\texttt{C}}
\noteblock{2}{\cellcolor{OliveGreen!40}}{\texttt{y}}{1}
\note{1}{\texttt{z}}
\noteunder{1}{\texttt{char}}
\emptyblock{1}{\cellcolor{Magenta!40}}
\emptyblock{1}{}
\]
В этом случае адрес базового объекта \verb"A" будет совпадать с адресом базового объекта \verb"C".
Сейчас можно добиться того же самого эффекта и с полем класса, только его надо пометить аттрибутом \verb"[[no_unique_address]] ".
Если предыдущее решение написать на языке вертикальной композиции (что это такое обсуждалось тут~\ref{}), то будет
\begin{cppcode}
struct D {
  B b;
  [[no_unique_address]] A a;
  C c;
  char z;
};
\end{cppcode}
У этого класса будет точно такой же layout как и у предыдущего.

\subsubsection{Наличие виртуальных методов}

В этом случае компилятор добавляет дополнительную информацию в структуру класса.
Куда и как он ее добавляет стандарт не специфицирует.
Однако, все современные компиляторы имплементируют виртуальные методы с помощью таблицы виртуальных функций (согласно спецификации Itanium C++ ABI).
И в классе хранится указатель на эту таблицу (чаще всего в начале).
А это значит, что у такого класса выравнивание будет максимальным $8$.
А размер вычисляется по той же самой схеме.
Например
\begin{cppcode}
struct A {
  virtual void f() const {}
  int x;
};

struct B {
  virtual void g() const {}
  short y;
};

struct C : A, B {
  virtual void h() const {}
  char z;
};
\end{cppcode}
В этом случае расположение в памяти выглядит так:
\[
\noteunder{16}{\texttt{A}}
\noteblock{8}{\cellcolor{Yellow!40}}{\texttt{vtable}}{2}
\noteblock{4}{\cellcolor{Orange!40}}{\texttt{x}}{1}
\noteblock{4}{}{\texttt{pad}}{2}
\noteunder{16}{\texttt{B}}
\noteblock{8}{\cellcolor{Yellow!40}}{\texttt{vtable}}{2}
\noteblock{2}{\cellcolor{OliveGreen!40}}{\texttt{y}}{1}
\noteblock{6}{}{\texttt{pad}}{2}
\noteunder{1}{\texttt{char}}
\noteblock{1}{\cellcolor{Magenta!40}}{\texttt{z}}{1}
\noteblock{7}{}{\texttt{pad}}{2}
\]
Обратите внимание, что новый виртуальный метод \verb"h" добавляется в виртуальную таблицу самого первого базового класса (в данном случае \verb"A").

Что будет, если \verb"A" не имел таблицы виртуальных методов, а \verb"B" имел?
\begin{cppcode}
struct A {
  int x;
};

struct B {
  virtual void g() const {}
  short y;
};

struct C : A, B {
  virtual void h() const {}
  char z;
};
\end{cppcode}
Тогда компилятор может переупорядочить базовые классы и считать, что у нас задан класс
\begin{cppcode}
struct C : B, A {
  virtual void h() const {}
  char z;
};
\end{cppcode}
В этом случае расположение данных в памяти будет
\[
\noteunder{16}{\texttt{B}}
\noteblock{8}{\cellcolor{Yellow!40}}{\texttt{vtable}}{2}
\noteblock{2}{\cellcolor{OliveGreen!40}}{\texttt{y}}{1}
\noteblock{6}{}{\texttt{pad}}{2}
\noteunder{4}{\texttt{A}}
\noteblock{4}{\cellcolor{Orange!40}}{\texttt{x}}{1}
\noteunder{1}{\texttt{char}}
\noteblock{1}{\cellcolor{Magenta!40}}{\texttt{z}}{1}
\noteblock{3}{}{\texttt{pad}}{2}
\]
Обратите внимание, что тут члены класса \verb"C" складываются по правилам выравнивания, как будто нет таблицы виртуальных функций.
Размер будет $24$ байта.

Теперь предположим, что ни один базовый класс не содержал виртуальных методов.
\begin{cppcode}
struct A {
  int x;
};

struct B {
  short y;
};

struct C : A, B {
  virtual void h() const {}
  char z;
};
\end{cppcode}
В этом случае заводится указатель на таблицу виртуальных методов и он скалывается в начале класса до помещения данных классов \verb"A" и \verb"B".
Расположение в памяти будет следующее:
\[
\noteblock{8}{\cellcolor{Yellow!40}}{\texttt{vtable}}{2}
\noteunder{4}{\texttt{A}}
\noteblock{4}{\cellcolor{Orange!40}}{\texttt{x}}{1}
\noteunder{2}{\texttt{B}}
\noteblock{2}{\cellcolor{OliveGreen!40}}{\texttt{y}}{1}
\noteunder{1}{\texttt{char}}
\noteblock{1}{\cellcolor{Magenta!40}}{\texttt{z}}{1}
\note{2}{\texttt{pad}}
\emptyblock{1}{}
\]
В этом случае размер будет $16$ байт.

\subsubsection{Наличие виртуальных базовых классов}

Ну и последний отвратительный пример
\begin{cppcode}
struct A {
  virtual void f() const {}
  int x;
};

struct B : virtual A {
  virtual g() const {}
  short y;
};

struct C : virtual A {
  virtual h() const {}
  float z;
};

struct D : B, C {
  virtual w() const {}
  char f;
};
\end{cppcode}
Выравнивание такого класса будет $8$, так как он содержит указатели.
В таком случае компилятор вырезает из \verb"B" и \verb"C" общую часть \verb"A" и оставляет только уникальную для них часть.
После чего общая часть \verb"A" складывается в конец класса.
Вот как будут выглядеть данные класса в памяти:
\[
\noteunder{16}{\texttt{B - A}}
\noteblock{8}{\cellcolor{Yellow!40}}{\texttt{vtable}}{2}
\noteblock{2}{\cellcolor{OliveGreen!40}}{\texttt{y}}{1}
\noteblock{6}{}{\texttt{pad}}{2}
\noteunder{12}{\texttt{C - A}}
\noteblock{8}{\cellcolor{Yellow!40}}{\texttt{vtable}}{2}
\noteblock{4}{\cellcolor{Violet!40}}{\texttt{z}}{1}
\noteunder{1}{\texttt{char}}
\noteblock{1}{\cellcolor{Magenta!40}}{\texttt{f}}{1}
\noteblock{3}{}{\texttt{pad}}{2}
\noteunder{16}{\texttt{A}}
\noteblock{8}{\cellcolor{Yellow!40}}{\texttt{vtable}}{2}
\noteblock{4}{\cellcolor{Orange!40}}{\texttt{x}}{1}
\noteblock{4}{}{\texttt{pad}}{2}
\]
В этом случае новые виртуальные методы из \verb"D" также помещаются в первый \verb"vtable", в данном случае это таблица для \verb"B".
Не спрашивайте как это работает в общем случае, самое важно, что мы получили из всех примеров выше -- по адресу класса найти адрес базового класса -- нетривиальная задача.

\subsection{Storage duration}
\label{section::StorageDuration}

Под \verb"storage" подразумевается память, выделенная под объект.
Таким образом тут я хочу обсудить время жизни памяти хранящей объект.
В языке есть следующие $4$ вида \verb"storage duration":
\begin{enumerate}
\item \verb"automatic".

\item \verb"static".

\item \verb"dynamic".

\item \verb"thread-local".
\end{enumerate}
Чтобы понять что это такое и что это значит, нам для начала надо понять какие виды памяти бывают.
Мы уже обсудили, что вся память для программиста выглядит как большой массив из байт.
Однако, операционная система делит ее на сегменты с разными правилами.
В зависимости от того, где лежат данные, у вас будет разное поведение по продолжительности жизни области памяти.

\subsubsection{Виды памяти}

Вся память используется как для хранения машинных инструкций (исполняемый код), так и для хранения данных.
Инструкции всегда хранятся в доступной только на чтение памяти.
У нас нет особенно возможности как-то повлиять на то, как в памяти лежат инструкции, а потому про эту часть особенно нечего сказать и мы сосредоточим свое внимание на памяти, которая хранит данные.

Можно выделить следующие сегменты в памяти:%
\footnote{Тут я по-хорошему должен разделять процессы и \verb"thread"-ы.
Но я буду неаккуратен в этом вопросе.}%
\footnote{Надо иметь в виду, что в зависимости от архитектуры железа и операционной системы, расположение указанных блоков в памяти может сильно отличаться.
Для нас это не очень принципиально, просто не зацикливайтесь именно на таком порядке данных.
Кроме того наличие зазоров между сегментами памяти не обязательно, например классикой считается ситуация, когда куча начинается сразу после глобальной статической памяти, но есть современные операционные системы, где это не так.}
\begin{enumerate}
\item \verb"static" память, которую можно разделить на следующие виды:
\begin{enumerate}
\item Read-only память для хранения машинных инструкций.

\item Read-only память для хранения константных данных вроде строковых литералов \verb$"123"$.

\item Глобальная static память программы.%
\footnote{Если вам это о чем-то говорит, я отношу сюда две секции \verb".data" и \verb".bss".}

\item Read-only глобальная static память под template-ы глобальных данных для \verb"thread_local" переменных.
Что это такое будет обсуждаться чуть ниже.%
\footnote{Взаимное расположение глобальных шаблонов для потоков и глобальной статической памяти может быть другим и зависит от конкретной имплементации.}
\end{enumerate}

\item Куча (Heap%
\footnote{Данный термин не имеет никакого отношения к структуре данных.
Это просто ходовое название, призванное показать, что в области кучи аллоцированные сегменты памяти лежат хаотично.
На практике там используется несколько разных структур данных.
}%
) для динамической аллокации памяти с помощью оператора \verb"new" или \verb"malloc" (и других методов аллокации).
Куча растет в сторону увеличения адресов в памяти.
Внутри кучи аллоцированная и деаллоцированная память находятся вперемешку.

\item Memory Mapping Segment (MMAP) сегмент.
В нем лежат:
\begin{enumerate}
\item Динамические библиотеки (\verb".so"/\verb".dll")

\item Для каждого \verb"thread"-а копия \verb"thread_local" template-ов.
\end{enumerate}

\item Для каждого \verb"thread"-а стек вызовов.
Стек вызова растет в сторону уменьшения адресов.
Его размер физически ограничен для всех \verb"thread"-ов единой константой зависящей от операционной системы.

\item Ядро операционной системы.
\end{enumerate}
Вот как можно визуализировать глобальную систему разных видов памяти:
\[
\datablock{4}{\cellcolor{Orange!40}}{\texttt{Global}}{\texttt{0x0000000000000000}}{1}
\middleblock{}
\noteunder{4}{\texttt{Heap}}
\emptyblock{4}{\cellcolor{Green!40}}
\middleblock{}
\noteunder{4}{\texttt{MMAP}}
\emptyblock{4}{\cellcolor{TealBlue!40}}
\middleblock{}
\noteunder{4}{\texttt{Stacks}}
\emptyblock{4}{\cellcolor{Salmon!40}}
\middleblock{}
\noteunder{4}{\texttt{OS}}
\emptyblock{4}{\cellcolor{BrickRed!40}}
\note{2}{\texttt{0xFFFFFFFFFFFFFFFF}}
\]

\subsubsection{Глобальная память}

Если смотреть только на глобальную память, то в соответствии с описанием выше у нас есть следующие зоны:
\[
\noteunder{16}{\texttt{Read-Only}}
\noteblock{8}{\cellcolor{Tan!40}}{\texttt{.text}}{2}
\noteblock{8}{\cellcolor{Violet!40}}{\texttt{.rodata}}{4}
\datablock{8}{\cellcolor{RedOrange!40}}{\texttt{Global Data}}{\texttt{.data/.bss}}{4}
\datablock{8}{\cellcolor{YellowGreen!40}}{\texttt{Thread Local Templates}}{\texttt{.tdata/.tbss}}{4}
\]
Давайте явно проговорим, что тут лежит.
\begin{enumerate}
\item В сегменте \verb".text" лежит исполняемый код.
Эта область памяти read-only.
То есть нельзя в run-time поменять исполняемый код модифицируя его бинарное представление в памяти.
Это сделано специально.

\item Во втором сегменте \verb".rodata" хранятся константные глобальные данные.
Вот что может тут лежать:
\begin{enumerate}
\item Строковые литералы:
\begin{cppcode}[numbers = none]
const char* str = "123";
\end{cppcode}
Тогда бинарное представление строки \verb$"123"$ лежит в этой части данных.
Стоит сказать, что указатель \verb"str" не лежит в этой области данных.
Он живет либо на стеке, либо на куче, либо в \verb".data".

\item Глобальные \verb"const" переменные (или статические \verb"const" переменные классов и методов, которые по сути тоже глобальные переменные но с ограниченной областью видимости).
Переменная обязана быть либо примитивным типом (встроенный целочисленный или с плавающей запятой) либо aggregate.%
\footnote{Смотри ниже~\ref{}.}
\begin{cppcode}[numbers = none]
struct A {
  static const int x = 0;
};

void f() {
  static const A a = {};
}

const int y = 42;

int main() {
}
\end{cppcode}
Переменные \verb"A::x", \verb"f::a"%
\footnote{Это условное обозначение.}%
, и глобальная \verb"y" лежат в константной глобальной памяти.

\item  Глобальные \verb"constexpr" и \verb"constinit" переменные (или статические \verb"constexpr" и \verb"constinit" переменные классов и методов).
\begin{cppcode}[numbers = none]
struct A {
  static constexpr int x = 0;
};

void f() {
  static constexpr A a = {};
}

constexpr int y = 42;

int main() {
}
\end{cppcode}

\item Таблица переходов для директивы \verb"switch".
Компилятор может имплементировать выбор в \verb"switch" в виде массива с указанием меток, куда надо перейти в коде.
В результате данный массив сохраняется в \verb".rodata".

\item Таблицы виртуальных функций.
Например класс
\begin{cppcode}[numbers = none]
struct A {
  virtual f() const {}
};
\end{cppcode}
хранит внутри себя указатель на таблицу виртуальных функций.
В данном случае таблица содержит указатель на метод \verb"A::f".
Данная таблица располагается в \verb".rodata" сегменте памяти.

\item Данные для поддержки RTTI.
Для обеспечения работы механизмов вроде \verb"dynamic_cast" и \verb"typeid".

\item Константы используемые в ассемблерных вставках.
Даже не хочу это комментировать.
\end{enumerate}

\item Третий сегмент <<Глобальных данных>> состоит из двух частей:
\begin{enumerate}
\item \verb".data"
Здесь лежат глобальные переменные (или статические переменные классов и функций), которые инициализируются конкретным значением во время компиляции.
\begin{cppcode}[numbers = none]
void f() {
  static int x = 42;
}

struct A {
  static int y;
};

int A::y = 42;

int z = 42;

int main() {
  return 0;
}
\end{cppcode}

\item \verb".bss".
Здесь лежат неинициализированные или инциализированные нулем глобальные переменные (или статические переменные классов и функций).
\begin{cppcode}[numbers = none]
void f() {
  static int x = 0;
}

struct A {
  static int y;
};

int A::y = 0;

int z;

int main() {
  return 0;
}
\end{cppcode}
\end{enumerate}

\item Четвертый сегмент Thread Local Segment Templates (TLS Templates) посвящен \verb"thread_local" данным.
Этот сегмент read-only и содержит мастер копии глобальных (и/или статических) переменных которые имеют \verb"thread_local" storage duration.
Эта память используется для дальнейшей инициализации \verb"thread_local" переменных для каждого конкретного \verb"thread"-а.
Он тоже состоит из двух частей.
\begin{enumerate}
\item \verb".tdata"
Здесь лежат глобальные переменные (или статические переменные классов и функций), которые инициализируются конкретным значением во время компиляции.
\begin{cppcode}[numbers = none]
void f() {
  thread_local int x = 42; // automatically static
}

struct A {
  thread_local static int y;
};

thread_local int A::y = 42;

thread_local int z = 42;

int main() {
  return 0;
}
\end{cppcode}
Переменные \verb"f::x", \verb"A::y" и \verb"z" будут расположены в этом сегменте, инициализированы указанными значениями и будут read-only.

\item \verb".tbss".
Здесь лежат неинициализированные или инциализированные нулем глобальные переменные (или статические переменные классов и функций).
\begin{cppcode}[numbers = none]
void f() {
  thread_local int x = 0; // automatically static
}

struct A {
  thread_local static int y;
};

thread_local int A::y = 0;

thread_local int z;

int main() {
  return 0;
}
\end{cppcode}
Переменные \verb"f::x", \verb"A::y" и \verb"z" будут расположены в этом сегменте, инициализированы нулями и будут read-only.
\end{enumerate}
Обратите внимание, что можно думать про этот раздел, что он содержит не сами \verb"thread_local" переменные, а лишь их изначальные значения (values).
Для каждого \verb"thread"-а надо заводить копию соответствующей переменной.
То есть нет единой общей переменной для всех потоков.
Так что в глобальной памяти хранятся лишь инициализаторы.
\end{enumerate}
Если кратко, то в статической части хранится глобальная информация.
К ней относятся: исполняемый код, глобальные переменные (константные и нет) и инициализаторы для \verb"thread_local" переменных.
Все глобальные данные будут лежать ту.

\subsubsection{Куча}

Данный сегмент памяти предназначен для динамического выделения памяти.
Следующие операции выделяют/удаляют память на куче:%
\footnote{Смотри так же раздел~\ref{section::MemoryAllocation}.}
\begin{enumerate}
\item оператор \verb"new" (кроме случаев вычисления во время компиляции).
Удаление с помощью оператор \verb"delete".

\item оператор \verb"new[]" (кроме случаев вычисления во время компиляции).
Удаление с помощью оператор \verb"delete[]".

\item \verb"std::malloc".
Удаление с помощью функции \verb"std::free".

\item \verb"std::calloc".
Удаление с помощью функции \verb"std::free".

\item \verb"std::realloc".
Удаление с помощью функции \verb"std::free".

\item \verb"std::aligned_alloc".
Удаление с помощью функции \verb"std::free".
\end{enumerate}
На практике операторы \verb"new" и \verb"new[]" внутри все так же обращаются к \verb"std::malloc" и компании.

Самое важное что надо понимать, что задача кучи -- выделять и освобождать сегменты памяти заранее неизвестного размера и делать это за разумное время.
Понятно, что при таких запросах память сильно фрагментируется и нужно хранить информацию о структуре выделенной памяти.
Представлять себе кучу стоит как-то так:
\[
\hspace{-0.5cm}
\noteunder{15}{\substack{\texttt{Global Memory}\\\texttt{.data/.bss}}}
\firstblock{}
\noteblock{8}{\cellcolor{Goldenrod!40}}{\texttt{bookkeeping}}{4}
\middleblock{}
\noteunder{32}{\texttt{Heap}}
\noteblock{8}{\cellcolor{Goldenrod!40}}{\texttt{bookkeeping}}{4}
\noteblock{2}{\cellcolor{Orange!40}}{\texttt{used}}{2}
\noteblock{2}{\cellcolor{Orange!40}}{\texttt{used}}{2}
\noteblock{4}{}{\texttt{free}}{2}
\noteblock{2}{\cellcolor{Orange!40}}{\texttt{used}}{2}
\noteblock{8}{\cellcolor{Orange!40}}{\texttt{used}}{2}
\noteblock{2}{}{\texttt{free}}{2}
\noteblock{4}{\cellcolor{Orange!40}}{\texttt{used}}{2}
\lastblock{}
\]
Здесь есть служебная информация $\begin{array}{|c|}\hline{\cellcolor{Goldenrod!40}}\\\hline\end{array}$, которая поддерживает состояние кучи, выделенная память $\begin{array}{|c|}\hline{\cellcolor{Orange!40}}\\\hline\end{array}$ и свободная память $\begin{array}{|c|}\hline{}\\\hline\end{array}$.

Служебная информация хранится в двух частях:
\begin{enumerate}
\item Глобальные данные для управления доступной памятью.
Такая информация хранится в \verb".data/.bss" сегменте.

\item Локальная информация о выделенных размерах сегментов.
Хранится локально внутри самой кучи.
Где именно и как -- зависит от имплементации.
Может храниться рядом с данными, а может в едином централизованном хранилище аллокатора (но в области памяти самой кучи).
\end{enumerate}
Поддержка такой информации -- дорогое удовольствие.
\begin{enumerate}
\item Эта информация не локализована в одном месте.
А потому необходимо прыгать по памяти, что ухудшает работу кэша.
В целом железо хорошо оптимизировано именно на последовательное чтение или запись данных.

\item Для поддержания служебной информации нужно выполнять дополнительный код, который занимает время.
\end{enumerate}
К плюсам кучи относится то, что она с точки зрения программиста -- бесконечный источник памяти.

\subsubsection{Стэк вызовов}

Давайте напомню (или сообщу) что логически реальную работу выполняет \verb"thread".
Именно он подключается на ядро процессора и именно он выполняет переданную на него функцию в качестве \verb"main".
И во время работы \verb"thread" может вызывать различные функции.
У каждой функции есть аргументы, возвращаемое значение и какие-то локальные переменные.
При выходе из функции аргументы и локальные переменные надо удалить, а возвращаемое значение положить по нужному адресу.
После чего вернуть исполнение вызвавшей функции в нужную точку.
Для поддержки подобного поведения и необходим стэк вызовов.
Более того, важно понимать, что стэк поддерживается на уровне железа.
У этого есть и плюсы и минусы.
\begin{itemize}
\item Главный плюс -- стек работает очень быстро.
\begin{enumerate}
\item В нем все данные лежат рядом.

\item Минимальная служебная информация для поддержания состояния стека.

\item Быстрая загрузка и выгрузка данных в специальные регистры.

\item Специальные машинные инструкции для стека.

\item Кэш оптимизирован под работу со стеком.
Железу легко предсказать какие данные понадобятся.
\end{enumerate}

\item Еще один плюс аппаратной поддержки стэка -- дополнительная защита корректности работы стека на аппаратном уровне.

\item Обратите внимание, что за счет поддержки на аппаратном уровне вызова функций, вы не можете использовать все возможные инструкции.
Например, если вы прыгните во время исполнения из одной функции в другую, то такая процедура не имеет смысла, если определено понятие вызова функции.
А так как железо поддерживает это понятие, то подобные трюки могут поставить железо в некорректное состояние.
В частности из-за этого в современных языках \verb"goto" не поддерживает возможность прыжка между функциями.%
\footnote{Так же полезно понимать, что когда Дейкстра писал о том, что \verb"goto" -- это зло, железо и языки программирования были совсем другими.
Тогда не было поддержки стека на уровне железа, потому что не было понятия вызова функции.
Тогда не было привычных нам контрольных элементов языка вроде \verb"if/else" или циклов \verb"while/for".
Все что было доступно программисту -- либо выполнить команду, либо прыгнуть с помощью \verb"goto" в любую точку программы.
Потому современное \verb"goto" это цветочки.}%

\item Так же стоит отметить, что размер стека сильно ограничен.
Даже внутри одной операционной системы размеры стеков для разных потоков могут быть разными.
Но обычно для операционной системы есть дефолтные значения размеров.
И конечно же у каждой операционной системы оно свое.
Стек съедается двумя вещами: локальными данными функций (включая аргументы и возвращаемое значение) и рекурсией (точнее ее глубиной).
\end{itemize}

Когда операционная система запускает процесс с программой, то создается некоторая экосистема для обеспечения функционирования процесса.
В эту экосистему как минимум входят:
\begin{enumerate}
\item Стэк вызовов.
В современных desktop компьютерах у каждого процесса есть свой основной стек вызова, который создается в $4$-ой зоне со стеками.
Если же вы запускаете \verb"thread"-ы из основного потока приложения, то их стэки вызовов уже создаются в TLS сегменте.

\item Thread Local Segment (TLS) для \verb"thread_local" данных.
Тут лежат \verb"thread_local" переменные для каждого стека.
Они инициализируются из мастер шаблона в \verb".tdata/.tbss" памяти.
\end{enumerate}
Так как в операционной системе может быть создано много процессов (а у каждого процесса может быть запущено много потоков), то для каждого из них уж точно нужно создать стэк вызовов и возможно еще и TLS сегмент.
Основные стэки, стеки потоков и TLS сегменты управляются по принципу работы кучи (heap).

\paragraph{Пример работы стэка}

Давайте разберемся на примере как работает стэк.
Я опущу всякие технические детали, которые не нужны для понимания, а нужны лишь для аккуратной имплементации работы со стэком на машинном языке.
Рассмотрим следующий код:
\begin{cppcode}
int f(int x, int y) {
  int w = x + y;
  w = 2 * w;
  return w;
}

int main() {
  int z = 0;
  z = f(2, 3);
  return 0;
}
\end{cppcode}
И давайте представим себе, что никаких RVO/NRVO оптимизаций не происходит.
Поговорим о том, как идейно работает стэк в самом наивном виде.
А уже после этого скажем пару слов о том, какие приблизительно оптимизации выполняются дополнительно.

Правила работы со стеком следующие.
Если мы заходим в функцию, то в начале выполняется предварительная работа:
\begin{enumerate}
\item В начале на стек складываются аргументы функции (их возможно перед этим надо вычислить и для этого понадобится вызов других функций).
В каком порядке они складываются стандарт не определяет.
И в $32$ битных системах были разные подходы.
В $64$ битных системах обычно складывают от последнего аргумента к первому.

\item складывается адрес, куда вернуть значение по \verb"return".
\end{enumerate}
Вот теперь можно вызывать функцию.
В этом месте логически открывается фрейм вызываемой функции.
На стек дополнительно пишется служебная информация.
Сразу скажем, что вся служебная информация состоит из двух указателей:
\begin{enumerate}
\item Один указывает на вершину стека, чтобы понимать где он заканчивается.

\item А другой указывает на начало фрейма функции, чтобы было удобно считать отступы относительно него.
\end{enumerate}
При открытии нового фрейма надо сохранить где лежит предыдущий адрес нового фрейма, чтобы суметь вернуться.
Его то и надо сохранить.
Внутри фрейма функции на стек складываются все локальные переменные функции.
При выходе из функции по сохраненному адресу возвращается значение по \verb"return" (если оно не \verb"void") и все локальные переменные уничтожаются вызовом деструкторов.
После чего фрейм очищается, что в реальности означает сдвиг двух указателей.

Теперь покажем это на примере выше.
\begin{cppcode}
int f(int x, int y) {
  int w = x + y;
  w = 2 * w;
  return w;
}

int main() {
  int z = 0;
  z = f(2, 3);
  return 0;
}
\end{cppcode}
Я буду рассматривать состояние стека с момента входа в функцию \verb"main".
Все данные на стек надо укладывать в соответствии с требованию по выравниванию.
Для простоты картины, я буду намерено игнорировать выравнивание, чтобы картинка не стала слишком большой.
При вызове $7$~строчки кода состояние стека будет следующим.%
{ % stack example macro scope
\newcommand{\args}{\cellcolor{Orange!40}}
\newcommand{\mainlocal}{\cellcolor{Emerald!40}}
\newcommand{\flocal}{\cellcolor{SeaGreen!40}}
\newcommand{\return}{\cellcolor{Mulberry!40}}
\[
\phantom{
\datablock{4}{\flocal}{\texttt{int}}{\texttt{w}}{1}
\datablock{8}{\return}{\texttt{int*}}{\texttt{\&z}}{1}
\datablock{4}{\args}{\texttt{int}}{\texttt{x}}{1}
\datablock{4}{\args}{\texttt{int}}{\texttt{y}}{1}
\datablock{4}{\mainlocal}{\texttt{int}}{\texttt{z}}{1}
}
\noteunder{4}{\substack{\texttt{OS specific}\\\texttt{calls}}}
\middleblock{}
\note{3}{\texttt{stack}\atop\texttt{bottom}}
\]
Теперь на строке~$8$ на стек складывается локальная переменная \verb"z".
\[
\phantom{
\datablock{4}{\flocal}{\texttt{int}}{\texttt{w}}{1}
\datablock{8}{\return}{\texttt{int*}}{\texttt{\&z}}{1}
\datablock{4}{\args}{\texttt{int}}{\texttt{x}}{1}
\datablock{4}{\args}{\texttt{int}}{\texttt{y}}{1}
}
\datablock{4}{\mainlocal}{\texttt{int}}{\texttt{z}}{1}
\noteleft{2}{\texttt{main}\atop\texttt{frame}}
\noteunder{4}{\substack{\texttt{OS specific}\\\texttt{calls}}}
\middleblock{}
\note{3}{\texttt{stack}\atop\texttt{bottom}}
\]
После чего надо вызвать функцию \verb"f" в строке~$9$.
Перед этим на стек надо сложить аргументы.
Будем складывать их от последнего к первому.
 \verb"z".
\[
\phantom{
\datablock{4}{\flocal}{\texttt{int}}{\texttt{w}}{1}
\datablock{8}{\return}{\texttt{int*}}{\texttt{\&z}}{1}
}
\datablock{4}{\args}{\texttt{int}}{\texttt{x}}{1}
\datablock{4}{\args}{\texttt{int}}{\texttt{y}}{1}
\datablock{4}{\mainlocal}{\texttt{int}}{\texttt{z}}{1}
\noteleft{2}{\texttt{main}\atop\texttt{frame}}
\noteunder{4}{\substack{\texttt{OS specific}\\\texttt{calls}}}
\middleblock{}
\note{3}{\texttt{stack}\atop\texttt{bottom}}
\]
Теперь  на стек надо положить адрес \verb"z" для того чтобы по этому адресу вернуть значение из функции.
\[
\phantom{
\datablock{4}{\flocal}{\texttt{int}}{\texttt{w}}{1}
}
\datablock{8}{\return}{\texttt{int*}}{\texttt{\&z}}{1}
\datablock{4}{\args}{\texttt{int}}{\texttt{x}}{1}
\datablock{4}{\args}{\texttt{int}}{\texttt{y}}{1}
\datablock{4}{\mainlocal}{\texttt{int}}{\texttt{z}}{1}
\noteleft{2}{\texttt{main}\atop\texttt{frame}}
\noteunder{4}{\substack{\texttt{OS specific}\\\texttt{calls}}}
\middleblock{}
\note{3}{\texttt{stack}\atop\texttt{bottom}}
\]
Теперь создаем локальную переменную \verb"w".
\[
\datablock{4}{\flocal}{\texttt{int}}{\texttt{w}}{1}
\noteleft{2}{\texttt{f }\atop\texttt{frame  }}
\datablock{8}{\return}{\texttt{int*}}{\texttt{  \&z}}{1}
\datablock{4}{\args}{\texttt{int}}{\texttt{x}}{1}
\datablock{4}{\args}{\texttt{int}}{\texttt{y}}{1}
\datablock{4}{\mainlocal}{\texttt{int}}{\texttt{z}}{1}
\noteleft{2}{\texttt{main}\atop\texttt{frame}}
\noteunder{4}{\substack{\texttt{OS specific}\\\texttt{calls}}}
\middleblock{}
\note{3}{\texttt{stack}\atop\texttt{bottom}}
\]
Теперь управление находится на строке~$2$ и в \verb"w" записывается значение выражения \verb"x + y".
При этом размер стека теперь не меняется.
Далее в строке~$3$ обновляем значение \verb"w".
В строке~$4$ нужно вернуть значение \verb"w" по адресу записанному на стеке.
Выполняются инструкции для записи результата в \verb"z".
Теперь нужно выйти из функции \verb"f".
В начале надо удалить все локальные переменные: \verb"w", \verb"&z", \verb"x", \verb"y".
Так как это \verb"int" и указатель \verb"int*", то деструкторы вызывать не надо.
Теперь мы просто выбрасываем все эти переменные со стека и получаем:
\[
\phantom{
\datablock{4}{\flocal}{\texttt{int}}{\texttt{w}}{1}
\datablock{8}{\return}{\texttt{int*}}{\texttt{\&z}}{1}
\datablock{4}{\args}{\texttt{int}}{\texttt{x}}{1}
\datablock{4}{\args}{\texttt{int}}{\texttt{y}}{1}
}
\datablock{4}{\mainlocal}{\texttt{int}}{\texttt{z}}{1}
\noteleft{2}{\texttt{main}\atop\texttt{frame}}
\noteunder{4}{\substack{\texttt{OS specific}\\\texttt{calls}}}
\middleblock{}
\note{3}{\texttt{stack}\atop\texttt{bottom}}
\]
Данное положение стека будет после выполнения строки~$9$.
Теперь в строке~$10$ выполняется запись $0$ по адресу определенного перед \verb"main" на стеке (куда в операционную систему вернуть значение).
И теперь в строке~$11$ происходит чистка стека.
В данном случае просто уменьшается сам стек.
\[
\phantom{
\datablock{4}{\flocal}{\texttt{int}}{\texttt{w}}{1}
\datablock{8}{\return}{\texttt{int*}}{\texttt{\&z}}{1}
\datablock{4}{\args}{\texttt{int}}{\texttt{x}}{1}
\datablock{4}{\args}{\texttt{int}}{\texttt{y}}{1}
\datablock{4}{\mainlocal}{\texttt{int}}{\texttt{z}}{1}
}
\noteunder{4}{\substack{\texttt{OS specific}\\\texttt{calls}}}
\middleblock{}
\note{3}{\texttt{stack}\atop\texttt{bottom}}
\]
Вот и все.
} % stack example macro scope

\paragraph{Дополнительные детали}

\begin{enumerate}
\item При вызове функции вида
\begin{cppcode}[numbers = none]
int f(int, int);
\end{cppcode}
Компилятор может изменить ее сигнатуру на 
\begin{cppcode}[numbers = none]
void f(int*, int, int);
\end{cppcode}
где первый аргумент -- указатель, который хранит адрес, куда нужно вернуть значение из функции.
Этот адрес может трактоваться как один из аргументов, а может подчиняться отдельным правилам.
Все зависит от конкретного ABI.

\item Каждый аргумент меньше $8$ байт занимает целиком $8$ байт и находится в левой (где адрес меньше) части $8$ байт (то есть padding-ги добавляются справа).
Например для функции
\begin{cppcode}[numbers = none]
void f(char a, bool b, int c);
\end{cppcode}
Будет следующее расположение:
\[
\newcommand{\Int}{\cellcolor{Orange!40}}
\newcommand{\Char}{\cellcolor{Magenta!40}}
\newcommand{\Bool}{\cellcolor{TealBlue!40}}
\firstblock{}
\noteunder{8}{\texttt{arg1}}
\noteleft{3}{\texttt{frame f  }}
\noteblock{1}{\Char}{\texttt{  a}}{1}
\emptyblock{7}{}
\noteunder{8}{\texttt{arg2}}
\noteblock{1}{\Bool}{\texttt{b}}{1}
\emptyblock{7}{}
\noteunder{8}{\texttt{arg3}}
\noteblock{4}{\Int}{\texttt{c}}{1}
\emptyblock{4}{}
\]

\item Первые несколько аргументов при вызове функции складываются сразу в регистры, а не на стек.
Количество аргументов, которые не складываются на стек зависит от архитектуры и операционной системы, например $4$ для x86 $64$ битной Windows и $6$ для $64$ битные Linux/macOS (System V Amd64).
Например на Windows вызов функции
\begin{cppcode}[numbers = none]
void f(int a, int b, int c, int d, int e, int f);
\end{cppcode}
отправит первые четыре аргумента \verb"a", \verb"b", \verb"c", \verb"d" в регистры напрямую, а \verb"e" и \verb"f" уйдут на стек.%
{\newcommand{\Int}{\cellcolor{Orange!40}}
\begin{align*}
\noteunder{8}{\texttt{register1}}
\noteblock{4}{\Int}{\texttt{a}}{1}
\emptyblock{4}{}
\nothingblock{1}
\noteunder{8}{\texttt{register2}}
\noteblock{4}{\Int}{\texttt{b}}{1}
\emptyblock{4}{}
\nothingblock{1}
\noteunder{8}{\texttt{register3}}
\noteblock{4}{\Int}{\texttt{c}}{1}
\emptyblock{4}{}
\nothingblock{1}
\noteunder{8}{\texttt{register4}}
\noteblock{4}{\Int}{\texttt{d}}{1}
\emptyblock{4}{}&
\\
\firstblock{}
\noteleft{3}{\texttt{frame f  }}
\noteunder{8}{\texttt{arg5}}
\noteblock{4}{\Int}{\texttt{  e}}{1}
\emptyblock{4}{}
\noteunder{8}{\texttt{arg6}}
\noteblock{4}{\Int}{\texttt{f}}{1}
\emptyblock{4}{}&
\end{align*}}
Имена конкретных регистров зависят от деталей имплементации.


\item Большие объекты (Windows -- больше $8$ байт, Linux -- больше $16$ байт) складываются внутри фрейма вызывающей функции, а в качестве аргумента кладется его адрес (в регистр или на стек).
Например пусть дана функция:
\begin{cppcode}[numbers = none]
struct A {
  double x[2];
};
void f(int x, A a, int y, int z, A b, bool w);
\end{cppcode}
Тогда состояние регистров и стека перед вызовом функции \verb"f" будет приблизительно такое:
\[
\newcommand{\Int}{\cellcolor{Orange!40}}
\newcommand{\Ptr}{\cellcolor{Orchid!40}}
\newcommand{\Bool}{\cellcolor{TealBlue!40}}
\newcommand{\Double}{\cellcolor{GreenYellow!40}}
\hspace{-1.5cm}\begin{aligned}
\noteunder{8}{\texttt{register1}}
\noteblock{4}{\Int}{\texttt{x}}{1}
\emptyblock{4}{}&
\nothingblock{1}
\noteunder{8}{\texttt{register2}}
\noteblock{8}{\Ptr}{\texttt{\&a}}{1}
\nothingblock{1}
\noteunder{8}{\texttt{register3}}
\noteblock{4}{\Int}{\texttt{y}}{1}
\emptyblock{4}{}
\nothingblock{1}
\noteunder{8}{\texttt{register4}}
\noteblock{4}{\Int}{\texttt{z}}{1}
\emptyblock{4}{}
\\
\firstblock{}
\noteleft{3}{\texttt{frame f  }}
\noteunder{8}{\texttt{arg5}}
\noteblock{8}{\Ptr}{\texttt{  \&b}}{1}
\noteunder{8}{\texttt{arg6}}
\noteblock{1}{\Bool}{\texttt{w}}{1}
\emptyblock{7}{}&
\datablock{16}{\Double}{\texttt{A}}{\texttt{a}}{1}
\datablock{16}{\Double}{\texttt{A}}{\texttt{b}}{1}
\end{aligned}
\]

\item Скалярный тип (так же POD типы малого размера) возвращаемый через \verb"return" и который может уместиться в $8$ байтах возвращается через регистр.
\begin{cppcode}
int f(int x, int y);
\end{cppcode}
Тогда
\[
\newcommand{\Int}{\cellcolor{Orange!40}}
\newcommand{\Ptr}{\cellcolor{Orchid!40}}
\newcommand{\Bool}{\cellcolor{TealBlue!40}}
\newcommand{\Double}{\cellcolor{GreenYellow!40}}
\newcommand{\return}{\cellcolor{Mulberry!40}}
\hspace{-1.5cm}\begin{aligned}
\noteunder{8}{\texttt{register0}}
\noteblock{4}{\return}{\texttt{return}}{2}
\emptyblock{4}{}
\nothingblock{1}
\noteunder{8}{\texttt{register1}}
\noteblock{4}{\Int}{\texttt{x}}{1}
\emptyblock{4}{}
\nothingblock{1}
\noteunder{8}{\texttt{register2}}
\noteblock{4}{\Int}{\texttt{y}}{1}
\emptyblock{4}{}&
\\
\firstblock{}&
\noteleft{3}{\texttt{frame f}}
\middleblock{}
\note{3}{\texttt{stack}\atop\texttt{bottom}}
\end{aligned}
\]

\item При вызове функции \verb"f" локальная переменная \verb"w" возвращается по \verb"return" и результат складывается в \verb"z".
Сейчас в стандарте языка явно прописаны правила для RVO/NRVO (раздел~\ref{}).
Правило такое, если мы знаем, какую переменную мы вернем из функции, то давайте не будем заводить под нее память, а используем адрес, куда ее надо вернуть, в качестве самой переменной.
То есть при вызове \verb"f" вместо \verb"w" будет использовать \verb"z" напрямую через ее адрес, который лежит на стеке (на самом деле в регистре).
Вернувшись к нашему примеру
\begin{cppcode}
int f(int x, int y) {
  int w = x + y;
  w = 2 * w;
  return w;
}

int main() {
  int z = 0;
  z = f(2, 3);
  return 0;
}
\end{cppcode}
Получим следующее представление регистров и стека перед вызовом \verb"f" приблизительно таким:
\[
\newcommand{\args}{\cellcolor{Orange!40}}
\newcommand{\mainlocal}{\cellcolor{Emerald!40}}
\newcommand{\flocal}{\cellcolor{SeaGreen!40}}
\newcommand{\return}{\cellcolor{Mulberry!40}}
\begin{aligned}
\noteunder{8}{\texttt{register0}}
\noteblock{4}{\return}{\texttt{return}}{2}
\emptyblock{4}{}
\nothingblock{1}
\noteunder{8}{\texttt{register1}}
\noteblock{4}{\args}{\texttt{x}}{1}
\emptyblock{4}{}
\nothingblock{1}
\noteunder{8}{\texttt{register2}}
\noteblock{4}{\args}{\texttt{y}}{1}
\emptyblock{4}{}&
\\
\firstblock{}
\noteleft{2}{\texttt{f }\atop\texttt{frame  }}
\emptyblock{4}{}
\datablock{4}{\mainlocal}{\texttt{int}}{\texttt{z}}{1}
\noteleft{2}{\texttt{main}\atop\texttt{frame}}
\noteunder{4}{\substack{\texttt{OS specific}\\\texttt{calls}}}
\middleblock{}
\note{3}{\texttt{stack}\atop\texttt{bottom}}
\end{aligned}
\]
Здесь я специально отобразил зазор между \verb"z" и началом фрейма для \verb"f", чтобы подчеркнуть, что фреймы выровнены по $16$ байтовым границам ($128$ бит).
Это нужно для корректного выравнивания данных для SIMD инструкций.%
\footnote{Single Instruction Multiple Data -- специальные машинные инструкции, которые могут сделать одну операцию над несколькими данными разом, например, сложить покомпонентно два массива из целых чисел.
Для этого данные массивы должны помещаться в специальные SIMD регистры.
Для таких данных обычно требуется особое выравнивание в памяти.}

\item На $64$ битных системах граница стека всегда выравнивается по $16$ байтам перед вызовом функции.
Если это условие не выполнено, то добавляется padding.
\end{enumerate}

\subsubsection{Memory Mapping Segment (MMAP)}

Еще один кусок памяти, который предназначен для выделения больших кусков памяти выровненных по размеру страниц памяти.%
% TO DO
\footnote{Видимо надо что-то сказать про страницы.}
Что может лежать в той области памяти:
\begin{enumerate}
\item Стеки вызовов всех \verb"thread"-ов, которые создает программа обычно выделяются тут.
То есть \verb"thread"-ы, где выполняется \verb"main" каждой программы имеет основной стек в куче из стеков, а вот второстепенные стеки уже выделяются тут.
\begin{cppcode}
#include <thread>

void function () {
  int x = 0;
}

int main() {
  std::thread t(function);
  t.join();
  return 0;
}
\end{cppcode}
В данном коде в $8$-ой строке кода создается \verb"thread" \verb"t", на котором запускается функция \verb"function".
Под вызовы функций создается стек в MMAP.
И локальная переменная \verb"x" из функции \verb"function" складывается на стек \verb"thread"-а \verb"t".
Так что переменная \verb"x" сама находится в MMAP сегменте.
\begin{cppcode}
#include <thread>

thread_local int s = 0;

void function() {
  s = 1;
}

int main() {
  s = 2;
  std::thread t(function);
  t.join();
  return 0;
}
\end{cppcode}


\item Активные \verb"thread_local" переменные, которые были инициализированы значениями из \verb".tdata/.tbss".

\item Динамические библиотеки.

\item Бинарник исполняемой программы может помещаться в эту часть.

\item Файлы целиком загруженные в память.

\item Большие аллокации из кучи могут перенаправляться сюда.

\item Разделяемая память между процессами операционной системы.%
\footnote{Видимо все таки надо разделять процессы и потоки.}

\item Специальные регистры железа могут непосредственно отображаться на сегмент памяти в этой зоне.
Изменение такого участка памяти напрямую изменяет соответствующие регистры в железе.
И в обратную сторону, соответствующее железо может напрямую писать в память обходя CPU.
\end{enumerate}
Для нас этот участок памяти больше интересен в связи с \verb"thread_local" переменными.
Соответствующий кусок памяти выделяется в момент создания \verb"thread"-а.

\subsubsection{Storage duration}

Теперь когда мы обсудили какой бывает память и как она выглядит, давайте обсудим, что означает \verb"storage duration" спецификатор в C++.
Всего есть $4$ опции:
\begin{enumerate}
\item \verb"automatic".

\item \verb"static".

\item \verb"dynamic".

\item \verb"thread_local".
\end{enumerate}
Для начала надо бы объяснить, а что такое \verb"storage".
Вообще говоря это логическая характеристика для компилятора, которую мысленно можно представлять как <<логически выделенная память>>.
Про \verb"storage" надо думать так: он показывает отмеченный сегмент в памяти, где объект находится.
Кроме того вместе со \verb"storage" идут правила, которые говорят в течение какого времени этот \verb"storage" вообще говоря существует в указанном сегменте памяти.
Важно понимать, что \verb"storage" это свойство объекта, а не области памяти.
Давайте рассмотрим два пример.
Первый:
\begin{cppcode}
struct A {
  A() {
    /*this makes A not to be implicit lifetime*/
  }
  ~A() {}
};

void* buffer = std::aligned_alloc(alignof(A), sizeof(A)); // allocate raw memory without objects
A* p = new (buffer) A(); // create an object of type A in the memory allocated before
                         // buffer provides a storage
p->~A();                 // destroy the object, the storage remains intact
std::free(buffer);       // deallocate the storage
\end{cppcode}
В этом примере можно думать про \verb"storage" как про сырую память выделяемую с помощью аллоцирующей функции \verb"std::aligned_alloc".
А объект уже создается внутри этого \verb"storage".
Вот другой пример
\begin{cppcode}
struct A {
  A() {
    /*this makes A not to be implicit lifetime*/
  }
  ~A() {}
};

int main() {
  alignas(A) unsigned char buffer[sizeof(A)];
  // create an array object buffer with automatic storage duration
  // so buffer has its own storage in memory on stack
  
  A* p = new (buffer) A();
  // create an object of type A in the storage provided by the object buffer
  // the object has dynamic storage duration
  
  p->~A();
  // destroy the object, the storage of A remains intact
  
} // buffer and its storage are destroyed here, hence the storage for A is destroyed here also
\end{cppcode}
Здесь же объект типа \verb"A" создается внутри объекта \verb"buffer".
Они имеют один и тот же размер и одно и то же выравнивание и занимают один и тот же кусок памяти.
Тем не менее, логически у них разные \verb"storage".
Более того, \verb"storage" для объекта типа \verb"A" находится внутри \verb"storage" для объекта \verb"buffer".
Оператор placement \verb"new" не аллоцирует, а переиспользует чужой \verb"storage" для размещения объекта.
Но с созданием объекта логически создается его \verb"storage".
При уничтожении объекта типа \verb"A" вызовом деструктора мы удаляем объект, но логически оставляем жить его \verb"storage", который ограничен временем жизни \verb"storage", в котором он находится, то есть \verb"storage" для \verb"buffer".
Потому только при выходе из \verb"main" уничтожается объект (массив) на стеке, и потом логически уничтожается его \verb"storage" при удалении этой памяти со стека.
Одна из причин почему мы разделяем \verb"storage" в этом примере -- \verb"buffer" и объект типа \verb"A" имеют разную storage duration в этом случае.
А один и тот же storage не может иметь две взаимоисключающие характеристики.

\paragraph{Automatic storage duration}

Следующие переменные имеют automatic storage duration:
\begin{enumerate}
\item Переменные принадлежащие block scope (ниже я поясню, что это такое) и не объявленные явно \verb"static", \verb"thread_local" или \verb"extern".
Время жизни памяти под такие объекты заканчивается при выходе из блока, в котором они были созданы.
\begin{cppcode}[numbers = none]
int main() {
  if (bool f = true; f) { // if scope begins and contains f
    int x;
  } // scope ends
  // f and x are not visible here
  { // scope begins
    int x; // completely new instance of x
  } scope ends
  return 0;
}
\end{cppcode}

\item Блок параметров функции (включает аргументы и тело функции).
Время жизни памяти под такие объекты заканчивается сразу после разрушения соответствующего объекта.
\begin{cppcode}[numbers = none]
int f(int a, int b) { // scope of f begins here and contains a and b
  int x = a + b;
  int w = 2 * x;
  return w - 1;
} // scope of f ends here, none of the variables are visible after this line
\end{cppcode}

\item Временные объекты так же имеют automatic storage duration.
\begin{cppcode}[numbers = none]
int main() {
  Matrix A = ...;
  Vector x = ...;
  Vector b = ...;
  Vector y = A * x + b;
}
\end{cppcode}
В последнем выражении вычисления проходят следующим образом
\begin{cppcode}[numbers = none]
  Vector tmp = A * x;
  Vector y = tmp + b;
\end{cppcode}
Временный объект, который я обозначил именем \verb"tmp", так же имеет automatic storage duration.
Его время жизни ограничено временем жизни полного выражения.
Полное выражение обычно означает строчка кода целиком, содержащая код порождающий временный объект.
Более точное описание будет ниже.
\end{enumerate}
Самое важное в этом виде storage duration -- вам не надо о нем задумываться, память под объекты выделяется и удаляется автоматически.

\paragraph{Block Scope}

Теперь давайте разбираться, что такое block scope.
Грубо говоря -- это код между фигурных скобок в исполняемом коде.
Вот что относится к block scope формально:
% TO DO
% Примеры
\begin{enumerate}
\item Блоки \verb"if" и \verb"switch".
\begin{cppcode}[numbers = none]
if (bool cond = true; cond) { // begins here and contains cond
  ...
} // ends here

switch (x) { // begins here
  case 0:
  break;
  default:
}// ends here
\end{cppcode}

\item Блоки циклов \verb"for", range-\verb"for", \verb"while" и \verb"do-while".
\begin{cppcode}[numbers = none]
for (int x = 0; x < 10; ++x) { // begins here and contains x
  ...
} // ends here

for (const auto& elem : Container{1, 2, 3}) { // begins here,
                                              // contains Container{1, 2, 3}
                                              // and elem reference
  ...
} // ends here

while (f) { // begins here
  ...
} // ends here

do { // begins here
  ...
} while (f); // ends here
\end{cppcode}

\item Блок обработки исключений \verb"try" и \verb"catch".
\begin{cppcode}[numbers = none]
try { // try scope begins here
  ...
} catch (...) { // try scope ends here
  ...
}

try {
  ...
} catch (...) { // catch scope begins here
  ...
} // catch scope ends here
\end{cppcode}

\item Блок из сгруппированных команд в фигурные скобки.
\begin{cppcode}[numbers = none]
int main() {
  { // begins here
    int x = 0;
    int y = 2 * x;
  } // ends here
  { // begins here
    int x = 2;
    int y = 2 - x;
  } // ends here
}
\end{cppcode}
\end{enumerate}

\paragraph{Полное выражение}

%Полное выражение -- это строка кода до \verb";".
% Главное исключение -- lifetime extension для const T& и T&&.
% Проблемы с биндом на подъобъект возвращаемый по ссылке.
% Проблемы с ranged-for loop и временным объектом в инициализаторе до C++23

% TO DO
TO DO


\paragraph{Static storage duration}

Переменные, для которых выполнены все условия ниже, имеют static storage duration:
\begin{enumerate}
\item Переменная лежит в пространстве имен или была объявлена с ключевым словом \verb"static" или \verb"extern".

\item Переменная не имеет thread storage duration.
\end{enumerate}
Например следующие переменные имеют static storage duration:
\begin{cppcode}[numbers = none]
extern int s;

void f() {
  static int s = 0;
}

struct A {
  static int s;
};

int A::s = 0;

int z;

namespace ext {
int z = 0;
}

int main() {
  return 0;
}
\end{cppcode}
Время жизни для блока памяти под такие сущности длится в течение исполнения программы.
С момента запуска \verb"main" и до момента выхода из нее, можно не задумываться о времени жизни статических объектов.
До вызова функции \verb"main" (в этот момент срабатывают конструкторы глобальных объектов) и после ее завершения (в этот момент срабатывают деструкторы глобальных объектов) уже нужно быть аккуратным с обращением к глобальным ресурсам.
Про правильное использование глобальных ресурсов можно посмотреть раздел~\ref{section::Global}.

\paragraph{Dynamic storage duration}

Это единственный вид storage duration, который требует ручного управления.
При этом ручное управление может быть у выделенной памяти и объекта вместе, а может быть только у объекта (например создание объекта внутри буфера на стеке).
Объекты созданные с помощью перечисленных ниже методов имеют dynamic storage duration:
\begin{enumerate}
\item С помощью оператора new.
\begin{cppcode}[numbers = none]
struct A {};

int main() {
  A* obj = new A(); // *obj has dynamic storage duration
  ...
  delete obj;       // destroy the object and free the memory
  
  void* buffer = std::malloc(sizeof(A)); // allocate raw memory
  new (buffer) A();  // placement new inherits dynamic storage duration
  ...
  buffer->~A();      // destroy the object
  std::free(buffer); // deallocate the memory
  
  unsigned char buffer1[sizeof(A)]; // automatic storage duration
  new (buffer1) A();  // dynamic storage duration, it is no inherited from buffer1
  ...
  buffer1->~A();      // destroy the object in the buffer
  return 0;
}
\end{cppcode}

\item Неявное создание любым способом (implicit lifetime объекты обсуждаются в разделе~\ref{section::ImplicitLifetime}).
Память под такие объекты пересекается с уже выделенной памятью.
Здесь есть прикольный случай, если вы сделаете \verb"memcpy" в буффер на стеке, то вы создадите внутри буффера (который сам имеет automatic storage duration) объект, который будет иметь dynamic storage duration.
Но при этом время жизни этого dynamic storage будет ограничено временем жизни буффера (есть отдельное правило подчинения одной памяти другой).
Это супер контринтуитивно, но связано с тем, что для таких неявных переменных в коде нет никаких объявлений.
А потому компилятору надо их относить к динамическим.
\begin{cppcode}
struct A { // this is implicit-lifetime class
  int x;
  int y;
};

int main() {
  alignas(A) unsigned char buffer[sizeof(A)]; // automatic storage duration
  int arr[2] = {42, 69};
  std::memcpy(buffer, arr, sizeof(A));
 // std::memcpy implicitly creates an object of type A in the buffer
 // this explains why we can access the memory through the pointer to A below
  A* a = static_cast<A*>(buffer); // *a has dynamic storage duration!
  ...
  buffer1->~A();  // destroy the object in the buffer.
  // There is not need, but we can end the lifetime explicitly
  return 0;
}
\end{cppcode}
Я обожаю этот пример.
Обратите внимание, что в качестве носителя памяти у нас выступает массив \verb"buffer" на стеке подходящего размера и правильно выровненный.
Сам массив и его элементы \verb"unsigned char" имеют automatic storage duration.
Однако, когда мы используем копирование с помощью \verb"std::memcpy", то эта функция создает в указанной области памяти любой объект с implicit-lifetime.
А потому мы можем обратиться в эту память по указателю на \verb"A".
Кроме того, из-за implicit-lifetime правила языка гласят, что объект типа \verb"A" размещенный в \verb"buffer" имеет dynamic storage duration.
Этот пример показывает, что объект носитель памяти \verb"buffer" и storage в определении storage duration для объекта \verb"A" -- это два разных понятия.
С точки зрения языка \verb"buffer" предоставил свой storage для storage-а переменной типа \verb"A".
И в этом случае время жизни storage для \verb"*a" будет не больше времени жизни storage для \verb"buffer".
То есть несмотря на dynamic storage duration объект все равно умрет по выходу из ближайшего scope.
Связано такое поведение с тем, что storage -- это понятие нужное компилятору, а не железу.
И переменная типа \verb"A" возникла внутри \verb"buffer" но при этом нигде нет кода где синтаксически выражено создание именно объекта класса \verb"A".
Создание происходит в строке~$9$ (об этом ниже после следующего примера), но компилятор это узнает по косвенным причинам в строке~$12$.
Изменение типа в строке~$12$ изменит создаваемый объект в строке~$9$.
Изменим код следующим образом:
\begin{cppcode}
struct A {
  int x;
  int y;
};

struct B {
  int x;
  int y;
};

int main() {
  alignas(A) unsigned char buffer[sizeof(A)];
  int arr[2] = {42, 69};
  std::memcpy(buffer, arr, sizeof(A));
  A* a = static_cast<A*>(buffer);
  B* b = static_cast<B*>(buffer); // this is still fine unless you access the data
  b->x = 1; // this is UB
  ...
  return 0;
}
\end{cppcode}
В строке~$14$ \verb"std::memcpy" может создать объект с implicit-lifetime.
В строке~$15$ компилятор понимает, что \verb"std::memcpy" действительно создала там объект типа \verb"A".
И значит возвращаясь к строке~$14$ компилятор считает, что функция \verb"std::memcpy" действительно неявно создала там объект типа \verb"A".
Теперь в строке~$16$ мы пытаемся обратиться к объекту класса \verb"A" через указатель на объект класса \verb"B", что нарушает правила работы с системой типов языка.
Мы не можем стучаться по чужому указателю в объект.
Если вам надо сменить объект в буффере, то для пользователей C++23 доступен следующий методв
\begin{cppcode}
int main() {
  alignas(A) unsigned char buffer[sizeof(A)];
  int arr[2] = {42, 69};
  std::memcpy(buffer, arr, sizeof(A));
  A* a = static_cast<A*>(buffer);
  B* b = std::start_lifetime_as<B>(buffer); // ends lifetime of A and starts lifetime of B
  b->x = 1; // OK
  ...
  return 0;
}
\end{cppcode}
В этом случае поля \verb"x" и \verb"y" в \verb"B" будут иметь те же самые значения.
Или если нет такой роскоши, то можно воспользоваться либо placement \verb"new"
\begin{cppcode}[numbers = none]
B* b = new (buffer) B;
\end{cppcode}
В этом случае поля \verb"x" и \verb"y" будут иметь неопределенные значения.
Другой способ -- использовать \verb"memmove"
\begin{cppcode}[numbers = none]
void* ptr = std::memmove(buffer, buffer, sizeof(B)); // ends lifetime of A
                                                     // and starts lifetime of B
B* b = static_cast<B*>(ptr);
\end{cppcode}
и надеяться, что компилятор это оптимизирует.
Обратите внимание, что я использую новый указатель, который возвращает placement \verb"new" или \verb"std::memmove".
Это связано с тем, что компилятор думает, что предыдущие указатели вроде \verb"buffer" или \verb"a" указывают на другой объект.
То есть компилятор каждому указателю логически приписывает то, на что он может ссылаться.
Подробнее об этом эффекте и как с этим всем работать будет подробнее рассказано в разделе~\ref{section::Launder}

\item Бросание исключение.
Память под такие объекты аллоцируется и деаллоцируется неспецифицированным способом.
\begin{cppcode}[numbers = none]
void f() {
  throw std::runtime_error("Ups"); // the exception thrown has dynamic storage duration
}
int main() {
  f();
  return 0;
}
\end{cppcode}
\end{enumerate}
Этот вид storage duration требует ручного управления памятью.
Конечно же это все надо делать с помощью RAII оберток (смотри также~\ref{}).

\paragraph{Thread local storage duration}

Все переменные объявленные с ключевым словом \verb"thread_local" имеют thread storage duration.
\begin{cppcode}
void f() {
  thread_local int x = 42; // automatically static
}

struct A {
  thread_local static int y;
};

thread_local int A::y = 42; // must repeat thread_local keyword

thread_local int z = 42;

int main() {
  return 0;
}
\end{cppcode}
Время жизни памяти для таких объектов длится в течение времени жизни соответствующего \verb"thread"-а, для которого они создавались.
Для каждого \verb"thread"-а существует своя уникальная копия такой переменной.
При использовании имени такой переменной используется переменная для данного \verb"thread"-а.
По сути эти переменные ведут себя как глобальные, просто из каждого потока видна своя уникальная версия этой переменной.
Нельзя обращаться к таким переменным до или после исполнения функции переданной на поток.
Точный момент инициализации таких переменных не специфицирован.
Они могут инициализироваться при первом обращении или до запуска потока.

\subsection{Lifetime}
\label{section::Lifetime}

%TO DO
Время жизни объекта в C++ является важным, потому что в языке у классов есть конструкторы и деструкторы, чья задача создать (или удалить) в нужном сегменте памяти объект в корректном состоянии.%
\footnote{В языке есть послабления для особых типов (об этом ниже), когда объекты могут создаваться неявно при обращении к сегменту памяти, но это в некоторой степени шаг в сторону совместимости с Си плюс гарантии, которые точно может дать компилятор.}
А потому обращаться к объекту в сегменте памяти, где он не был создан -- концепция противоречащая дизайну языка.
Уничтожать в методе объекта самого себя -- тоже не имеющая смысла вещь, хотя формально можно обвешаться костылями, чтобы это не было UB.
Тем не менее, понимание того, когда, где и как живут объекты, должно дать вам правильную ментальную модель.
И одна из важных привычек -- не делать бессмысленную хрень вроде \verb"delete this" и прочую ерунду.

\subsubsection{Начало lifetime}

Понятие lifetime (время жизни) объекта -- логическое понятие в языке, которое означает период времени, во время которого объект считается живым и к  нему логически можно обратиться.
В начале сформулируем правило, которое корректно работает во всех случаях кроме маленького количества исключений, а потом уже обозначим исключения явно.

\paragraph{Общее правило}

Время жизни объекта начинается, когда выполнены одновременно ОБА условия:
\begin{enumerate}
\item Выделена правильно выровненная память нужного размера.

\item Закончилась инициализация объекта (включая инициализацию по умолчанию или инициализацию тривиальным конструктором).
\end{enumerate}
Выделение памяти вам почти никогда не надо контролировать руками, единственным исключением являются объекты с dynamic storage duration.
А за создание отвечают в основном конструкторы (и еще некоторые функции работающие с сырой памятью могут создавать объекты неявно, об этом будет подробнее в разделе~\ref{section::ImplicitLifetime}).
Выделение памяти и объекта может быть синтаксически разделено (dynamic storage duration):
\begin{cppcode}[numbers = none]
struct A {
  int x;
};

int main() {
  /*memory operations*/
  void* p = std::aligned_alloc(alignof(A), sizeof(A)); // allocating a dynamic storage on heap
  /*object operations*/
  A* q = new (p) A; // create an A object default-initialized
  q->~A();          // destroy the object
  /*memory operations*/
  std::free(p);     // deallocate the memory reserved
  return 0;
}
\end{cppcode}
а может выполняться в виде одной операции (automatic, static, and thread-local storage duration):
\begin{cppcode}[numbers = none]
int x; // static storage duration
       // the memory is allocated by runtime
       // the value of the object is zero-initialized

struct A {
  int x;
};

int main() {
  int y; // automatic storage duration
         // the memory is allocated on stack by runtime
         // the object is initialized but the value is indeterminate
  return 0;
}
\end{cppcode}
Инициализация может проходить неявно для implicite-lifetime типов с помощью специальных функций работы с памятью:
\begin{cppcode}[numbers = none]
struct A {
  int x;
};

int main() {
  int y = 1;
  alignas(A) unsigned char buffer[sizeof(A)];
  std::memcpy(buffer, &y, sizeof(int)); // this implicitly creates an object of type A
                                        // in the buffer because of next line
  A* a = static_cast<A*>(buffer);
  return 0;
}
\end{cppcode}

\paragraph{Исключения}

Важно понимать, что инициализация объекта НЕ означает инициализацию подобъекта.
В основном такие проблемы крутятся вокруг нескольких примеров:
\begin{enumerate}
\item Поля в \verb"union" или подобъекты полей в \verb"union".
Для таких полей \verb"lifetime" начинается только если этот член \verb"union"-а был инициализирован
\begin{cppcode}[numbers = none]
union U {
  int x = 1;
  bool y;
};

U u; // the field u.x is active an alive, u.y is not alive
\end{cppcode}
или был сделан активным:
\begin{cppcode}[numbers = none]
union U {
  int x;
  bool y;
};

U u;     // neither u.x nor u.y is alive
u.x = 5; // u.x is alive
u.y = 6; // u.y is alive now
\end{cppcode}

\item Копирующие и перемещающие конструкторы и операторы присваивания для \verb"union" делают в копии активным поле, которое было активно в исходном \verb"union":
\begin{cppcode}[numbers = none]
union U {
  int x;
  bool y;
};

U u = {1}; // u.x is active an alive
U w = u;   // w.x is also active and alive
\end{cppcode}

\item Элементы массива не стартуют свой lifetime при создании массива с помощью \verb"std::allocator::allocate":
\begin{cppcode}[numbers = none]
struct A {
  int x;
  float y;
};

std::allocator<A> alloc;
A* p = alloc.allocate(10); // creates an array of type A[10]
                           // but its elements are not created
\end{cppcode}

\item Поля и базовы классы implicit-lifetime (что это такое в разделе~\ref{section::ImplicitLifetime}) типа, которые сами не implicit-lifetime.
\begin{cppcode}[numbers = none]
struct A { // implicit-lifetime
  int x;
  std::string y;
};

A a;
alignas(A) unsigned char buffer[sizeof(A)];

std::memcpy(buffer, &a, sizeof(A));
A* a = static_cast<A*>(buffer);

a->x = 1;     // OK
a->y = "123"; // UB, a->y is not alive yet
\end{cppcode}
\end{enumerate}

\subsubsection{Методы аллокации памяти}
\label{section::MemoryAllocation}

% TO DO
TO DO

\subsubsection{Явное создание объекта}

Существуют следующие способы явного создания объекта:
\begin{enumerate}
\item Определение (definition) объекта
\begin{cppcode}[numbers = none]
void f() {
  static int x; // a definition, starts lifetime
}

struct A {
  static int y; // a declaration, does not start lifetime
};

int A::y;       // a definition, starts lifetime

int z;          // a definition undergo zero-initialization, starts lifetime

extern int w;   // a declaration, does not start lifetime

int main() {
  int u;        // this is already a definition, starts lifetime
}
\end{cppcode}

\item Выражение \verb"new"
\begin{cppcode}[numbers = none]
A* p = new A(1, "123");   // starts lifetime of an object of type A

alignas(A) unsigned char buffer[sizeof(A)]; // memory buffer
new (buffer) A(1, "123"); // placement new starts lifetime of an object of type A
\end{cppcode}

\item Выражение \verb"throw"
\begin{cppcode}[numbers = none]
void f() {
  throw std::runtime_error("Ups"); // starts lifetime of the exeption thrown
}
\end{cppcode}

\item Изменение активного члена \verb"union".
\begin{cppcode}[numbers = none]
union A {
  int x;
  bool y;
};

A a;         // starts lifetime of the object of type A but none of the subobjects is alive
a.x = 1;     // starts the lifetime of a.x.
a.y = false; // ends the lifetime of a.x and starts the lifetime of a.y
\end{cppcode}
Определение переменной \verb"a" вообще говоря не стартует lifetime никакого из ее членов.
Дефолтный конструктор доступен только если все члены \verb"union" имеют тривиальные конструктор и деструктор (сгенерированные компилятором).
Если же мы используем класс с нетривиальным конструктором как \verb"std::string" (или инициализируем одно из полей явно), то дефолтного конструктора у нас нет и надо обязательно задать деструктор иначе он будет удален.
Это, правда, не сильно поможет в корректной работе с \verb"union" ибо надо еще знать кто там лежал, прежде чем удалять.
\begin{cppcode}[numbers = none]
union A {
  int x;
  bool y;
  std::string z;
  ~A() { /*need to know which member is active to delete*/ }
};

A a = {1}; // no default constructor because of the member z with non-trivial constructor
           // makes a.x active and starts its lifetime
new (&s.z) std::string("123"); // ends the lifetime of a.x and starts the lifetime of a.z
a.z.~basic_string(); // ends the lifetime of a.z
\end{cppcode}
Обратите внимание, что в случае нетривиального конструктора и деструктора, нам приходится явно вызывать placement \verb"new" и деструктор.
В случае типов с implicit-lifetime такая необходимость тоже есть, но можно активировать новый член инициализацией его.
\begin{cppcode}[numbers = none]
union A {
  int x;
  bool y;
  std::string z;
  ~A() { /*need to know which member is active to delete*/ }
};

A a = {1}; // no default constructor because of the member z with non-trivial constructor
           // makes a.x active and starts its lifetime
new (&s.z) std::string("123"); // ends the lifetime of a.x and starts the lifetime of a.z
a.x = 1;   // ends lifetime of a.z but does not free its resources and starts lifetime of a.x
\end{cppcode}
Есть еще всякие мутные правила про копирующие и перемещающие операции для \verb"union" в стандарте, но они все берутся из того, что \verb"union" не полноценный объект, он не знает всей информации о себе, чтобы банально обеспечить свою семантику.
Например для удаления или копирования \verb"union" надо знать кто в нем сейчас лежит, но сам \verb"union" этого не знает, а потому кто-то снаружи должен обеспечить семантику для класса.
И вокруг этого выростает огромное количество проблем и костылей.
В целом я не рекомендую ни в каком виде использовать \verb"union".
В STL есть \verb"std::variant", который является полноценным типом, который умеет менеджерить свои ресурсы.

\item Вычисление выражения с созданием временного объекта результата.
\begin{cppcode}[numbers = none]
int main() {
  Matrix A = ...;
  Vector x = ...;
  Vector b = ...;
  Vector y = A * x + b;
}
\end{cppcode}
В последнем выражении вычисления проходят следующим образом
\begin{cppcode}[numbers = none]
  Vector tmp = A * x;
  Vector y = tmp + b;
\end{cppcode}
В этом примере lifetime временного объекта, который я обозначил именем \verb"tmp", стартует после вычисления выражения \verb"A * x".
 
\item Вызов \verb"std::allocator::allocate" начинает lifetime массива объектов, но не его индивидуальных элементов (он никогда не стартует lifetime элементов даже для implicit lifetime~\ref{section::ImplicitLifetime}).
\begin{cppcode}[numbers = none]
struct A {
  int x;
  float y;
};

std::allocator<A> alloc;
A* p = alloc.allocate(10); // alloc creates an array A[10] and starts its lifetime,
                           // does not create elements
\end{cppcode}
\end{enumerate}

\paragraph{Несколько объектов по одному адресу}

По одному адресу можно создать 0 или более объектов.
Важно понимать, что расположение двух объектов по одному адресу может быть как корректным, так и не корректным с точки зрения языка.
Вот примеры корректных случаев:
\begin{enumerate}
\item Первый простой пример такой
\begin{cppcode}[numbers = none]
struct A {
  int x;
};

A a;
\end{cppcode}
Тогда объект \verb"a" и \verb"a.x" располагаются по одному адресу и занимают одно и то же место в памяти.
При этом \verb"storage" для \verb"a.x" логически находится внутри \verb"storage" для \verb"a" и равен ему по размеру.
Технически это два разных объекта.

\item Массив объектов
\begin{cppcode}[numbers = none]
struct A {
  int x;
  int y;
};

A arr[10];
\end{cppcode}
Здесь массив \verb"arr" и его первый элемент \verb"arr[0]" располагаются по одному адресу, но они занимают разную память.

\item Пример когда один объект предоставляет свой storage для другого объекта:
\begin{cppcode}[numbers = none]
alignas(A) unsigned char buffer[sizeof(A)];
A* p = new (buffer) A(1, "123");
\end{cppcode}
Здесь массив \verb"buffer" и \verb"*p" находятся по одному адресу и занимают одну и ту же память.

\item Пустые базовые классы за счет empty base optimization дают другой пример
\begin{cppcode}[numbers = none]
struct A {};
struct B {};
struct C : A, B {};
\end{cppcode}
Подъобекты \verb"A" и \verb"B" создаются по адресу \verb"C".

\item Аналогичная ситуация для полей структуры получается так:
\begin{cppcode}[numbers = none]
struct A {};
struct B {};
struct D {
  [[no_unique_address]] A a;
  [[no_unique_address]] B b;
};
\end{cppcode}
Поля \verb"D::a" и \verb"D::b" создаются по адресу объекта \verb"D".
\end{enumerate}

\paragraph{Пошаговый пример вызова конструктора}

Давайте покажем, как выглядит память во время вызова конструктора класса.
Рассмотрим следующие классы.
\begin{cppcode}
struct A {
  int x;
};

struct B {
  int x;
};

struct C : A {
  int x;
  B b;
};
\end{cppcode}
Теперь разберем исполнение первой строки в функции \verb"main" вот в таком примере:
\begin{cppcode}[numbers = none]
int main() {
  C c;
  return 0;
}
\end{cppcode}
В ней вызывается дефолтный конструктор для \verb"c".
Давайте сразу нарисуем layout для объекта типа \verb"C" в памяти:
\[
\newcommand{\Int}{\cellcolor{Orange!40}}
\datablock{4}{\Int}{\texttt{A}}{\texttt{A::x}}{2}
\datablock{4}{\Int}{\texttt{int}}{\texttt{x}}{1}
\datablock{4}{\Int}{\texttt{B}}{\texttt{b}}{1}
\]
Размер класса будет $12$ байт, выравнивание -- $4$ байта.
Теперь предположим мы зашли во фрейм функции \verb"main" на стеке:
\[
\nothingblock{12}{}
\mathmakebox[0cm][r]{\firstblock{}}
\noteleft{2}{\texttt{main}\atop\texttt{frame}}
\noteunder{4}{\substack{\texttt{OS specific}\\\texttt{calls}}}
\middleblock{}
\note{3}{\texttt{stack}\atop\texttt{bottom}}
\]
Будет вызываться дефолтный конструктор для \verb"C", который логически выглядит так:
\begin{cppcode}[numbers = none]
C::C() : A(), /*nothing for x, */ B() { }
\end{cppcode}
Перед созданием объекта класса в начале выделяется память под объект нужного размера и нужным образом выровненная.
\[
\newcommand{\Int}{\cellcolor{Orange!40}}
\noteunder{12}{\texttt{raw memory}}
\emptyblock{12}{}
\noteleft{2}{\texttt{main}\atop\texttt{frame}}
\noteunder{4}{\substack{\texttt{OS specific}\\\texttt{calls}}}
\middleblock{}
\note{3}{\texttt{stack}\atop\texttt{bottom}}
\]
После этого вызываются конструкторы для всех базовых классов.
В данном случае дефолтный конструктор для \verb"A".
Так как у поля \verb"A::x" нет инициализатора и оно встроенного типа, то значение его будет неопределенным после отработки конструктора.
\[
\newcommand{\Int}{\cellcolor{Orange!40}}
\datablock{4}{\Int}{\texttt{A}}{\texttt{A::x}}{2}
\noteunder{8}{\texttt{raw memory}}
\emptyblock{8}{}
\noteleft{2}{\texttt{main}\atop\texttt{frame}}
\noteunder{4}{\substack{\texttt{OS specific}\\\texttt{calls}}}
\middleblock{}
\note{3}{\texttt{stack}\atop\texttt{bottom}}
\]
После этого вызываются конструкторы для всех полей класс в порядке их объявления сверху вниз.
В начале поле \verb"x".
Так как оно встроенного типа \verb"int", то у него нет конструктора.
Так как нет инициализирующего значения, то это поле остается с неопределенным значением.
\[
\newcommand{\Int}{\cellcolor{Orange!40}}
\datablock{4}{\Int}{\texttt{A}}{\texttt{A::x}}{2}
\datablock{4}{\Int}{\texttt{int}}{\texttt{x}}{1}
\noteunder{4}{\texttt{raw memory}}
\emptyblock{4}{}
\noteleft{2}{\texttt{main}\atop\texttt{frame}}
\noteunder{4}{\substack{\texttt{OS specific}\\\texttt{calls}}}
\middleblock{}
\note{3}{\texttt{stack}\atop\texttt{bottom}}
\]
И теперь вызывается дефолтный конструктор для \verb"B".
\[
\newcommand{\Int}{\cellcolor{Orange!40}}
\datablock{4}{\Int}{\texttt{A}}{\texttt{A::x}}{2}
\datablock{4}{\Int}{\texttt{int}}{\texttt{x}}{1}
\datablock{4}{\Int}{\texttt{B}}{\texttt{b}}{1}
\noteleft{2}{\texttt{main}\atop\texttt{frame}}
\noteunder{4}{\substack{\texttt{OS specific}\\\texttt{calls}}}
\middleblock{}
\note{3}{\texttt{stack}\atop\texttt{bottom}}
\]

Теперь объект создан на стеке.


\subsubsection{Implicit lifetime}
\label{section::ImplicitLifetime}

Некоторые операции могут неявно создавать объекты в выделенной области памяти, начиная их \verb"lifetime".
Типы для которых такое возможно относятся к implicit-lifetime типам.
По определению объекты следующих типов являются implicit-lifetime:
\begin{enumerate}
\item Скалярные типы (возможно с const/volatile):
\begin{multicols}{2}
\begin{enumerate}
\item Арифметические типы:
\begin{enumerate}
\item Целые типы:
\begin{enumerate}
\item \verb"bool"

\item Character types:
\begin{itemize}
\item \verb"char", \verb"signed char", \verb"unsigned char".

\item \verb"char8_t" (начиная с C++20).

\item \verb"char16_t", \verb"char32_t", \verb"wchar_t".
\end{itemize}

\item Знаковые типы: 
\begin{itemize}
\item \verb"signed char", \verb"short", \verb"int", \verb"long", \verb"long long".
\end{itemize}


\item Беззнаковые типы:
\begin{itemize}
\item \verb"unsigned char", \verb"unsigned short", \verb"unsigned int", \\
\verb"unsigned long", \verb"unsigned long long".
\end{itemize}
\end{enumerate}

\item Типы с плавающей запятой:
\begin{itemize}
\item \verb"float", \verb"double", \verb"long double".
\end{itemize}
\end{enumerate}


\item enum-ы
\begin{enumerate}
\item unscoped, то есть \verb"enum".

\item scoped, то есть \verb"enum class".
\end{enumerate}

\item Указатели:
\begin{enumerate}
\item Указатель на объект.

\item Указатель на функцию.
\end{enumerate}

\item Указатель на член класса:
\begin{enumerate}
\item Указатель на данные класса.

\item Указатель на метод класса.
\end{enumerate}

\item \verb"std::nullptr_t".

\item \verb"std::meta::info" (since C++26)
\end{enumerate}
\end{multicols}

\item Классы, которые относятся к implicit-lifetime (возможно с const/volatile).
% Нужны примеры попадает в один класс, но не попадает в другой
Есть два типа таких классов:%
\footnote{Тут имеется в виду одно из условий, либо первое, либо второе.
То есть оба пункта определяют implicit-lifetime классы.}
\begin{enumerate}
\item Aggregate класс, у которого деструктор не является user-provided.
Вот полный перечень условия включая определение aggregate (выполнено одновременно):
\begin{enumerate}
\item Нет user-provided деструкторов.

\item Нет user-declared или унаследованных конструкторов (начиная с C++20)

\item Нет приватных или защищенных  не статических членов объявленных непосредственно в самом классе

\item Нет виртуальных базовых классов (начиная с C++17)

\item Нет приватных или защищенных прямых базовых классов  (начиная с C++17)

\item Нет виртуальных методов
\end{enumerate}
Вот пример
\begin{cppcode}[numbers = none]
TO DO
\end{cppcode}

\item Класс, для которого одновременно выполнены следующие два условия:
\begin{enumerate}
\item У класса есть как минимум один доступный тривиальный конструктор.
В данном случае доступный означает что он не удален и не заблокирован ограничениями вроде неподходящего типа аргументов, concept-ов или \verb"enable_if" шаблонов (и подобных механизмов).
То есть такой конструктор разрешено выбрать компилятору во время процедуры поиска метода.
При этом это может быть копирующий или перемещающий конструктор, главное чтобы он был тривиальный (сгенерирован компилятором).

\item Класс имеет не удаленный тривиальный деструктор.
\end{enumerate}
Вот пример
\begin{cppcode}[numbers = none]
TO DO
\end{cppcode}
\end{enumerate}

\item Массивы любых типов (возможно с const/volatile).
\end{enumerate}
Если грубо резюмировать все определение, то из встроенных типов мы запретили ссылки и функции, а класс должен либо быть aggregate, либо иметь тривиальные деструктор и конструктор (причем любой из).

\paragraph{Подобъекты implicit-lifetime не обязаны быть implicit-lifetime}

Обратите внимание, что если объект \verb"x" является implicit-lifetime, то его подъобъект, вообще говоря может не быть implicit-lifetime.
То есть даже если вы создали в памяти какой-то  implicit-lifetime объект, то это еще не значит, что ко всем его полям или базовым классам можно обращаться, потому что они еще не созданы.
Вот два важных примера.
\begin{enumerate}
\item Структура содержащая не implicit-lifetime объект.
\begin{cppcode}
struct A {
  int x;
  std::string y;
};

A* ptr = static_cast<A*>(std::aligned_alloc(alignof(A), sizeof(A)));

ptr->x = 42;                     // OK
ptr->y = "123";                  // UB! ptr->y does not exist here
new(&ptr->y) std::string("123"); // OK if we comment the previous line
                                 // Also you have to delete the field ptr->y manually
\end{cppcode}
В данном случае структура \verb"A" является implicit lifetime классом (есть тривиальный конструктор и тривиальный не удаленный деструктор).
С другой стороны \verb"std::string" не является implicit lifetime классом, потому что он менеджерит ресурс внутри себя, его конструктор и деструктор выполняют нетривиальную работу.
Потому при обращении к выделенной памяти с помощью \verb"std::malloc" по указателю на \verb"A", мы можем предполагать, что объект такого типа был неявно создан в этом сегменте памяти.
Но мы не можем полагать, что все его поля и базовые классы были созданы, например, поле \verb"y" не было создано, а значит и обращаться к нему нельзя.

\item Массив из не implicit-lifetime объектов.
\begin{cppcode}
using Strings = std::string[5];

Strings* strings = static_cast<Strings*>(
                       std::aligned_alloc(alignof(Strings), sizeof(Strings)));

strings[0] = "123"; // UB strings[0] does not exist here
new(&(strings[0])) std::string("123"); // OK
\end{cppcode}
Таким образом мы создаем объект \verb"strings" являющегося массивом, но элементы массива не создаются.
То есть важно понимать, что массив -- это отдельный тип и объект этого типа можно создать не умея создавать сами объекты.
\end{enumerate}

\subsubsection{Неявное создание объекта}

\paragraph{Что мы хотим и в чем проблема}

Так как C++ вырос поверх языка Си, то было бы хорошо, чтобы обращения с памятью, которые можно делать в Си, были бы законны и в C++.
Например, в Си можно любой кусок памяти правильного размера и правильно выровненный трактовать как некую структуру и обращаться корректно к ее полям.
Однако в отличие от Си, C++ построен по аксиоматическому принципу и содержит много логических сущностей вроде объектов, их \verb"storage", их \verb"lifetime" и прочее.
Рассмотрим следующий пример:
\begin{cppcode}[numbers = none]
void* p = std::malloc(10);
unsigned char* buffer = static_cast<unsigned char*>(p);
\end{cppcode}
Мы аллоцировали сырую память и вернули адрес начала на соответствующий сегмент памяти в указатель \verb"p".
Но этот указатель не имеет типа.
В следующей строке мы скастовали этот указатель уже в типизированный, то есть мы говорим, что хотим трактовать область памяти выделенную \verb"std::malloc" как массив из $10$ элементов типа \verb"unsigned char".
Но теперь, если мы обращаемся к элементам этого массива
\begin{cppcode}[numbers = none]
buffer[0] = 'c';
\end{cppcode}
возникает логическая проблема.
Операция \verb"static_cast" позволяет создать массив \verb"buffer" как объект, занимающий выделенную область памяти.
Однако, сами элементы массива в эту память никто не помещал.
Синтаксически мы нигде не провели инициализацию элементов типа \verb"unsigned char" в каждой ячейке памяти, как того требует определение для \verb"lifetime"-а (см~\ref{section::Lifetime}).
При этом нельзя просто разрешить обращение к памяти
\begin{cppcode}[numbers = none]
buffer[0] = 'c';
\end{cppcode}
трактовать моментом создания элементов.
Мы можем с таким же успехом обратиться к этой области памяти по указателю другого типа, что должно будет создать там новый объект.
\begin{cppcode}[numbers = none]
bool* q = static_cast<bool*>(p);
q[0] = true;
\end{cppcode}
И тогда ломается логика, первый элемент массива \verb"buffer" внезапно трактуется как \verb"bool", а остальные элементы как \verb"unsigned char".%
\footnote{
Это не единственная проблема.
Если посмотреть defect report \href{https://www.open-std.org/jtc1/sc22/wg21/docs/papers/2020/p0593r6.html}{P0593R6}, то возможность создавать типы в любом месте памяти при обращении к этой памяти с помощью указателя нужного типа влечет модель, которая совпадает с моделью \verb"effective type" в языке Си.
И с этой моделью есть куча проблем связанных с анализом обращения к памяти через типизированные указатели (не убить оптимизации и при этом хоть как-то помочь программисту с ошибками).
В самом языке Си до сих пор нет хорошего решения этой проблемы, что считается признаком того, что не надо идти по такому пути.
}
Вместо этого в C++ некоторые отдельные операции наделяются возможностью неявно создавать объекты некоторых типов.
Например, функция \verb"std::malloc" наделяется такой возможностью.
Как она понимает, что за объекты она создала?
Она как бы смотрит в будущее как будет впервые использоваться данная память.
И по этой информации компилятор понимает, что \verb"std::malloc" неявно создал именно такой-то тип.
Важно обратить внимание, что повторный трюк с той же областью памяти уже будет UB.
Рассмотрим следующий пример
\begin{cppcode}
struct A {
  char x;
  bool y;
};

struct B {
  char x;
  bool y;
};

int main() {
  void* buffer = std::malloc(sizeof(A) * 10);
  A* a = static_cast<A*>(buffer); // this line makes std::malloc to create an array
                                  // of 10 objects of type A
  B* b = static_cast<B*>(buffer); // this is still fine unless you access the data
  b[0]->x = 'c';                  // this is UB
  ...
  return 0;
}
\end{cppcode}
В данном примере первое использование буффера памяти в по указателю на объект типа \verb"A" заставляет компилятор решить, что \verb"std::malloc" (не сообщив никому) создала неявно внутри \verb"buffer" массив из структур \verb"A".
И начиная со строчки вызова \verb"std::malloc" компилятор считает, что по указателю \verb"buffer" или \verb"a" лежит массив из $10$ объектов типа $A$.
В строчке~$15$ мы просто просим трактовать указанный адрес как адрес на другую структуру.
Это само по себе не проблема, проблема, если мы теперь обратимся по этому адресу.
Теперь в строчке~$16$, где идет обращение к первому элементу, мы видим, что мы обращаемся к области памяти где сейчас находится живой объект типа \verb"A", но через указатель на тип \verb"B".
Однако, объекта типа \verb"B" в этой области памяти никто не создавал (даже неявно, потому что точка создания -- не точка обращения к памяти, а точка выделения памяти).
А обращение к объекту за пределами его \verb"lifetime" -- это UB.

\paragraph{Механизмы неявного создания}

Сразу скажем, что неявное создание объектов возможно только для implicit-lifetime типов.
Ниже перечислены все механизмы по неявному созданию объектов:
\begin{enumerate}
\item Если объект с implicit-lifetime будет создан внутри массива из \verb"unsigned char" или \verb"std::byte", то считается, что операция создающая сам массив так же неявно создает и объект.%
\footnote{\label{ref::constexpr}Кроме вычислений во время компиляции \verb"constexpr/consteval", в этом случае просто не создается никаких объектов.}
Например
\begin{cppcode}[numbers = none]
struct A {
  int x;
  int y;
};

alignas(A) unsigned char buffer[sizeof(A)]; // create an object of type A implicitly
                                            // because of the next line of code
A* ptr = static_cast<A*>(buffer); 
\end{cppcode}
Здесь важно подчеркнуть, что \verb"static_cast" не является точкой создания объекта типа \verb"A".
Эта строчка нужна компилятору, чтобы понять, что при конструировании массива, в нем неявно возник объект типа \verb"A".
Точкой создания будет момент конструирования самого массива \verb"buffer", который предоставляет свой \verb"storage" для хранения объекта типа \verb"A".

\item Вызов аллоцирующих функций (они неявно создают объекты внутри аллоцируемой памяти при необходимости):
\begin{enumerate}
\item operator \verb"new".%
\textsuperscript{\ref{ref::constexpr}}
\begin{cppcode}[numbers = none]
void* buffer = ::operator new (sizeof(int), std::align_val_t{alignof(int)});
int* p = static_cast<int*>(buffer); // because of this line
                                    // the operator new creates int implicitly
\end{cppcode}
Если бы мы никогда не обращались к этому участку памяти, то не было бы неявно создано ни одного объекта.
Именно этот оператор вызывается при использования синтаксиса
\begin{cppcode}[numbers = none]
int* p = new int;
\end{cppcode}

\item operator \verb"new[]".%
\textsuperscript{\ref{ref::constexpr}}
\begin{cppcode}[numbers = none]
void* buffer = ::operator new[] (sizeof(int), std::align_val_t{alignof(int)});
int* p = static_cast<int*>(buffer); // because of this line
                                    // the operator new creates int implicitly
\end{cppcode}
Именно этот оператор вызывается при использования синтаксиса
\begin{cppcode}[numbers = none]
int* p = new int[10];
\end{cppcode}

\item \verb"std::malloc".
\begin{cppcode}[numbers = none]
void* buffer = std::malloc(sizeof(char));
char* p = static_cast<char*>(buffer); // because of this line
                                      // std::malloc creates char implicitly
\end{cppcode}
Данная функция возвращает адрес выровненный так, чтобы уж точно все фундаментальные типы языка можно было в него поместить.
Так же надо учесть, что на современных x64 машинах стек выровнен по $16$ байтам.
В результате, на современных $64$ битных системах такая функция скорее всего возвращает адрес выровненный по $16$ байтам.

\item \verb"std::calloc".
\begin{cppcode}[numbers = none]
void* buffer = std::calloc(10, sizeof(char));
char* p = static_cast<char*>(buffer); // because of this line
                                      // std::calloc creates array of 10 char-s
\end{cppcode}
Эта функция эквивалентна вызову
\begin{cppcode}[numbers = none]
void* buffer = std::malloc(10 * sizeof(char));
\end{cppcode}
Требования по выравниванию такие же, как и у \verb"std::malloc".

\item \verb"std::realloc".
\begin{cppcode}[numbers = none]
void* buffer = std::malloc(10 * sizeof(char));
char* p = static_cast<char*>(buffer); // because of this line
                                      // std::calloc creates array of 10 char-
s
void* tmp = std::realloc(p, 12 * sizeof(char)); // reallocate new buffer with 12 char-s
char* q = static_cast<char*>(tmp);
\end{cppcode}
Здесь важное замечание, что функция \verb"std::realloc" неявно создает implicit-lifetime объекты, если таковые уже существовали в выделенной памяти.
В силу того, что мы уже обратились к памяти выделенной \verb"std::malloc" через указатель \verb"char*", то теперь в \verb"buffer" живет массив \verb"char" из $10$ элементов.
И потому при вызове \verb"realloc" создается массив длины $12$, где первые $10$ элементов точно должны быть \verb"char".
А значит и все элементы будут типа \verb"char".
Кроме того, последние два новых элемента будут тоже созданы и они будут в неопределенном состоянии (не инициализированы).

\item \verb"std::aligned_alloc".
\begin{cppcode}[numbers = none]
void* buffer = std::aligned_alloc(alignof(int), sizeof(int));
int* p = static_cast<int*>(buffer); // because of this line
                                    // std::alighed_alloc creates an int
\end{cppcode}
На практике эта функция нужно только для вызова аллокаций для типов с выравниванием больше $16$ байт (для x64 систем).
На практике для встроенных типов эта функция просто вызывает \verb"std::malloc", хотя это и не гарантируется.
Кроме того у этой функции второй аргумент должен делиться на первый, что дает ограничения, которых нет в \verb"std::malloc".
\end{enumerate}

\item Вызов следующих функций для копирования представления объекта (object representation), в таком случае эти объекты создаются в области памяти назначения:
\begin{enumerate}
\item \verb"std::memcpy".
\begin{cppcode}[numbers = none]
struct A {
  int x;
  int y;
};

A a;
alignas(A) unsigned char buffer[sizeof(A)];
A* p = std::memcpy(buffer, &a, sizeof(A)); // this creates an object of type A
                                           // in the buffer
p->x = 1;                                  // This is OK
\end{cppcode}
Здесь важное замечание, что функция \verb"std::memcpy" как и \verb"std::realloc" неявно создает implicit-lifetime объекты, если таковые уже существовали в выделенной памяти.
В силу того, что мы копируем содержимое объекта из переменной \verb"a", где располагается живой объект, то после выполнения \verb"memcpy" по адресу назначения \verb"buffer" автоматически создается объект типа \verb"A".
При этом даже нет необходимости обращаться к этому участку памяти.

\item \verb"std::memmove".
\begin{cppcode}[numbers = none]
struct A {
  int x;
  int y;
};

A a;
alignas(A) unsigned char buffer[sizeof(A)];
A* p = std::memmove(buffer, &a, sizeof(A)); // this creates an object of type A
                                            // in the buffer

p->x = 1;                                   // This is OK
\end{cppcode}
Данная функция подчиняется тем же правилам по созданию объектов, что и \verb"std::memcpy".

\item \verb"std::bit_cast".
Данная операция создает новый объект из битового представления старого объекта.
\begin{cppcode}[numbers = none]
double x = 1.69;
uint64_t y = std::bit_cast<uint64_t>(x); // creates uint64_t from
                                         // the bit representation of double
\end{cppcode}
\end{enumerate}

\item Вызов следующих функций, которые специально стартуют lifetime объекта в выделенной памяти (с++23):
\begin{enumerate}
\item \verb"std::start_lifetime_as"
\begin{cppcode}[numbers = none]
TO DO
\end{cppcode}

\item \verb"std::start_lifetime_as_array"
\begin{cppcode}[numbers = none]
TO DO
\end{cppcode}
\end{enumerate}
\end{enumerate}


\subsubsection{Завершение lifetime}

% TO DO
% Для каждого способа нужен пример.
% Нужно сказать, что есть завершение lifetime-а, а есть корректное завершение lifetime-а и это не одно и то же.
Время жизни объектов заканчивается в следующих случаях:
\begin{enumerate}
\item Если тип объекта не класс, то \verb"lifetime" заканчивается в момент разрушения объекта (возможно явным вызовом деструктора).
Обожаю это правило, потому что надо еще догадаться, как вызвать деструктор явно у скалярных типов.
Оказывается это можно сделать следующими способами:
\begin{cppcode}[\nonumbers]
template<class T>
void destroy(T* p) {
  p->~T(); // this works in templates only, direct call p->~int() is a compilation error
}

using Int = int;
int main() {
  int x;
  x.~Int(); // it works via using only x.~int() is a compilation error
  
  alignas(int) unsigned char buffer[sizeof(int)];
  int* p = new (buffer) int;
  destroy(p);
  p->~Int(); // also works
}
\end{cppcode}
На встроенных типах можно вызывать деструктор несколько раз на одном объекте, это не UB.
Такой вызов не дает никакого кода для встроенных типов.
Вызов деструктора через шаблон или \verb"using" связано с тем, чтобы библиотечный код мог чистить память от объектов для классов и встроенных типов одинаково.
Иначе было бы печально рассматривать всегда два случая: встроенный тип или класс, и у нас бы получилась джава вместо C++.

\item Если тип объекта -- класс, то \verb"lifetime" заканчивается в начале вызова деструктора.
\begin{cppcode}[\nonumbers]
struct A {};

int main() {
  A a;
  {
    A b;
  } // b.~A() is called automatically
  alignas(A) unsigned char buffer[sizeof(A)];
  A* p = new (buffer) A;
  p->~A(); // the object is destroyed manually
  return 0;
} // a.~A() is called automatically
\end{cppcode}

\item Память, которую занимал объект, освободили или она переиспользована другим объектом, который не расположен внутри текущего (не является nested object, про это замечание ниже).%
\footnote{Важно понимать разницу между <<завершить \verb"lifetime">> и <<корректно завершить \verb"lifetime">>.
Например если бы в этом примере вместо \verb"int" мы создавали \verb"std::string" внутри \verb"buffer", то при выходе из \verb"f", аллоцированная строкой память на куче утечет.
То есть язык позволяет нам завершать время жизни объекта даже, если это не корректно с точки зрения использования ресурсов.}
\begin{cppcode}[\nonumbers]
// storage is deleted
void f() {
  alignas(int) unsigned char buffer[sizeof(int)];
  int* p = new(buffer) int;
} // here buffer is deleted, hence the lifetime of int inside the buffer is terminated

// storage is reused
alignas(bool) unsigned char buffer[sizeof(bool)];
bool* p = new(buffer) bool;
char* q = new(buffer) char; // reuse the memory for char
                            // ends lifetime of bool and begins lifetime of char
\end{cppcode}
Если удаляется объект внутри переменной (локальной или статической), то надо до вызова ее деструктора туда поместить другой объект иначе это UB.
\begin{cppcode}[\nonumbers]
struct A {
  A() {/*non-trivial*/}
  ~A() {/*non-trifial*/}
};
int main() {
  A a;
  a.~A();
} // UB the destructor is called the second time for a
// need to add A* p = new(&a) A; for solving this issue
// but you have to access the new object via p, or use std::launder on &a in general case
\end{cppcode}
Так же обратите внимание, что язык не требует вызова деструктора для завершения lifetime объекта.
Это конечно же может повести к проблемам с выделенными ресурсами, но это ортогональная для Lifetime-а задача.
Важно, что в случае делегирования конструктора, мы не создаем несколько объектов друг за другом в одной области памяти.
Например, пусть у нас есть следующий класс:
\begin{cppcode}[\nonumbers]
struct A {
  A(int x) : x(x) {}
  A(int x, char y): A(x), y(y) {}
  
  int x;
  char y;
};

A a(1, 'c');
\end{cppcode}
В этом случае логически происходит следующее.
При создании переменной \verb"a" в начале выделяется память под объект.
Потом срабатывает конструктор \verb"A(x)", который инициализирует поле \verb"x".
И компилятор считает, что именно в этот момент стартовал \verb"lifetime" объекта \verb"a" (несмотря на то, что он еще не достроен).
После этого уже инициализируется поле \verb"y" и выполняется тело второго конструктора.
Если в теле второго конструктора бросить исключение, то обязательно сработает деструктор, из-за того, что нужно разрушить все уже созданные объекты.
А так как \verb"lifetime" у объекта \verb"a" уже стартовал, то надо его закончить.
Хотя снаружи это выглядит так, как будто мы сначала построили объект с помощью первого конструктора, а потом как в случае placement \verb"new" поверх достроили недостающую часть второго объекта.
\end{enumerate}

\paragraph{Nested objects}

К nested objects относятся следующие случаи:
\begin{enumerate}
\item не статическая переменная класса:
\begin{cppcode}[\nonumbers]
struct A {
  static int x; // NOT nested inside A
  int y;        // nested inside A
  int z;        // nested inside A
};
\end{cppcode}

\item Базовый класс:%
\footnote{Если относится к базовым классам как к нестатической переменной класса, у которой нет имени, то логика более или менее понятна.
Тут нет аналога статической переменной.}
\begin{cppcode}[\nonumbers]
struct A {
  int x;
};

struct B : A { // A is nested inside B
  int y;
};
\end{cppcode}

\item Элементы массива:
\begin{cppcode}[\nonumbers]
struct A {
  int x;
};

A arr[10]; // elements arr[0], arr[1], ... , arr[9] are nested in arr
\end{cppcode}

\item Поля \verb"union":
\begin{cppcode}[\nonumbers]
union A {
  int x;   // x is nested in A
  bool y;  // y is nested in A
  short z; // z is nested in A
};
\end{cppcode}
\end{enumerate}

\paragraph{Замечания}

Некоторые замечания по времени жизни объектов.
\begin{itemize}
\item Lifetime объекта находится внутри lifetime для \verb"storage" этого объекта.
\begin{cppcode}[\nonumbers]
void f() {
  alignas(int) unsigned char buffer[sizeof(int)];
  int* p = new(buffer) int;
} // here buffer is deleted, hence the lifetime of int inside the buffer is terminated
\end{cppcode}

\item Lifetime ссылки начинается при инициализации и заканчивается по правилам, как будто это scalar type объект (то есть как указатель).
\begin{cppcode}[\nonumbers]
int x;
{
  int& y = x; // start lifetime of y
}             // end lifetime of y
\end{cppcode}

\item Lifetime объекта, на который ссылается ссылка может закончиться раньше, чем lifetime ссылки, что ведет к проблеме висячей ссылки (dangling reference).
\begin{cppcode}[\nonumbers]
int& f() {
  int x;
  return x;
}

int& y = f();
\end{cppcode}
В этом случае мы возвращаем ссылку на локальную переменную, которую запоминаем в ссылке \verb"y".
Но после выхода из \verb"f" переменная \verb"x" уже закончила свой lifetime.
Тем не менее, есть правило продления жизни временным объектам, об этом ниже в разделе~\ref{section::TemporaryObjects}.

\item Lifetimes не статических членов класса и базовых классов начинается в порядке инициализации.
А заканчивается в обратном порядке к порядку инициализации.
\begin{cppcode}[\nonumbers]
struct A {
  int x;
  int y;
};

struct B : A {
  double z;
};

B b;
\end{cppcode}
В этом случае сначала инициализируется базовый класс \verb"A" (здесь иницализация идет в порядке: сначала \verb"x", потом \verb"y"), а потом поле \verb"z".
При разрушении выполняем все в обратном порядке.
В начале разрушаем \verb"z", а потом разрушаем базовую часть \verb"A".
При разрушении базовой части, разрушаем объекты в обратном порядке вначале \verb"y", потом \verb"x".
\end{itemize}

\subsubsection{Переиспользование storage}

% TO DO
TO DO

\subsubsection{Transparently replaceable объекты}

% TO DO
TO DO

\subsubsection{Зачем нужна \texttt{std::launder}}
\label{section::Launder}
% TO DO
TO DO

%\subsection{Тип}
%
%% TO DO
%TO DO

\subsection{Values (значения объектов)}
\label{section::ObjValue}

\subsubsection{Value representation}

% TO DO
TO DO

\subsubsection{Value Category}
\label{section::ValueCategory}

% TO DO
TO DO

\subsubsection{Временные объекты}
\label{section::TemporaryObjects}

% TO DO
TO DO

%\subsection{Имя}
%
%% TO DO
%TO DO

\subsection{Преобразование типов}

% TO DO
TO DO

\subsubsection{Преобразование значений (values)}

% TO DO
TO DO

\subsubsection{Преобразование ссылок и указателей}

% TO DO
TO DO

\subsubsection{Strict Aliasing Rule}

% TO DO
TO DO

\subsubsection{Reachable memory}

% TO DO
TO DO

\subsubsection{Pointer interconvertibility}

% TO DO
TO DO

\subsubsection{Виды \texttt{cast}-ов}
% TO DO
TO DO

\paragraph{\texttt{static\_cast}}
% TO DO
TO DO

\paragraph{\texttt{reinterpret\_cast}}
% TO DO
TO DO

\paragraph{\texttt{const\_cast}}
% TO DO
TO DO

\paragraph{\texttt{dynamic\_cast}}
% TO DO
TO DO

\paragraph{\texttt{std::bit\_cast}}
% TO DO
TO DO

\paragraph{Си-\texttt{cast}}
% TO DO
TO DO

\subsection{Имплементация holder-ов}

% TO DO
TO DO

\subsubsection{Владение без счетчика \texttt{std::optional}}
% TO DO
TO DO

\subsubsection{Владение со счетчиком}
% TO DO
TO DO

\subsubsection{Владение несколькими типами \texttt{std::variant}}
% TO DO
TO DO

\subsubsection{Специальный случай \texttt{std::expected}}
% TO DO
TO DO

\subsubsection{Small Object Optimization}
% TO DO
TO DO


\subsection{CPU caches}

% TO DO
TO DO

\subsection{Single Instruction Multiple Data (SIMD)}

% TO DO
TO DO

\end{document}
